\RequirePackage{fix-cm}
\documentclass[a4paper, twoside]{report}
%% Language and font encodings
\usepackage[english]{babel}
\usepackage[utf8x]{inputenc}
\usepackage[T1]{fontenc}
\usepackage{textcomp}
\usepackage{caption}
\usepackage{lipsum}
\usepackage{mathtools}
\usepackage{amsmath,amssymb,amsfonts}
\usepackage{rotating}
\usepackage{graphicx}
\usepackage{float}
\graphicspath{{Figures/}}
\usepackage{amsmath}
\usepackage{amssymb}
\usepackage{xcolor}
\usepackage{soul}
\usepackage{array}
\usepackage{url}
\usepackage{blindtext}
\usepackage{comment}
\usepackage{tabularx}
\usepackage{booktabs}
\usepackage{ragged2e}
\usepackage{totalcount}
\usepackage{setspace}

\newcommand{\automatedwordcount}{%
  \immediate\write18{cat introduction.tex problem.tex methods.tex | texcount -1 -sum -q - > word_count.sum}%
  \input{word_count.sum}%
}

\usepackage[numbers,sort&compress]{natbib}
\usepackage[hidelinks]{hyperref}
\usepackage{breakurl}

\def\UrlBreaks{\do\/\do-}

\includeonly[
%usepackage{hyperref}

\makeatletter
\newcount\SOUL@minus %% <- without this line the page number would be reset to 0
\makeatother

%\usepackage{fancyhdr}


% ------------------------------------------------------------------------------------------
% FORMATTING

\usepackage[compact]{titlesec}
\titlespacing{\section}{0pt}{2ex}{2ex}
\titlespacing{\subsection}{0pt}{1ex}{1ex}
\titlespacing{\subsubsection}{0pt}{0.5ex}{1ex}

\titleformat{\section}
  {\normalfont\fontsize{18}{40}\bfseries}{\thesection}{1em}{}
  
%% Sets page size and margins
\usepackage[a4paper,top=2.8cm,bottom=2.8cm,left=3cm,right=3cm,marginparwidth=1.75cm]{geometry}

%% Useful packages


% ------------------------------------------------------------------------------------------

% \usepackage[compact]{titlesec}
% \titlespacing{\section}{12pt}{0pt}{0pt}
% \titlespacing{\subsection}{12pt}{0pt}{0pt}
% \titlespacing{\subsubsection}{12pt}{0pt}{0pt}

% \titleformat{\section}
%   {\normalfont\fontsize{12}{24}\bfseries}{\thesection}{24pt}{\doublespacing}

% % Redefine the formatting for \subsubsection
% \titleformat{\subsubsection}[runin]{\normalfont\itshape}{\thesubsubsection}{1em}{}[:]

% \titleformat{\chapter}[display]
%   {\normalfont\fontsize{12}{12}\bfseries\centering} % Adjust the font size and line spacing as desired
%   {\chaptertitlename\ \thechapter}
%   {12pt} % Adjust the horizontal space between "Chapter #" and the title
%   {}

% %% Sets page size and margins
% \usepackage[letter,top=2.8cm,bottom=2.8cm,left=3cm,right=3cm,marginparwidth=1.75cm]{geometry}

% %% Useful packages
% \usepackage{amsmath}
% \usepackage{graphicx}
% \usepackage[colorlinks=true, allcolors=black]{hyperref}




% ------------------------------------------------------------------------------------------
%% Making a list of acronyms

\usepackage{acro}
\usepackage[acronym]{glossaries}
\makeglossaries 

% ------------------------------------------------------------------------------------------
%% Making a list of acronyms
\usepackage{acro}
\usepackage[acronym]{glossaries}
\makeglossaries

% Make your acronym list here:

\newacronym{ad}{AD}{Alzheimer's Disease} % The argument in the first pair of braces is the reference label, the second is the abbreviation in the text, and the third is the full meaning of the abbreviation.
\newacronym{who}{WHO}{World Health Organization} 


% ------------------------------------------------------------------------------------------
% NOTE: Use this section for the temporary manuscript only
%             Omit this section when printing the final, revised manuscript
%             You should also remove the line: "\linenumbers" below 
%             (before "\chapter{Introduction and Literature Review}

% =============================================================
% PART A — INTRODUCTION
% =============================================================

%% MARQUEZ START
\section{Background of the Study}
The public transportation system has long been a commodity to people; it operates on different terrains, has different forms, and has predetermined routes depending on their geographical location and the needs of commuters \cite{Tironi2018}. With an increase in urbanization over time, the demand for transportation has also increased. Despite the predicted decline in urbanization rates in Southeast Asia, the Philippines is experiencing the opposite \cite{UN_WUP2018}. Due to the increased urbanization and the continuous increase in population (especially in urban areas), people start looking for different alternatives in transportation \cite{Cervero2000}. With the increasing number of motor vehicles on the road, the traffic congestion worsens, which leads to commuters finding other means to quickly reach their destination. In Metro Manila, severe traffic congestion, stemming from insufficient road infrastructure and poor traffic mitigation policies, has resulted in a highly inefficient public transport system \cite{Spatio2026}. As a response, the Philippine government launched the EDSA Bus Carousel in 2020. It is the only busway system in Metro Manila implemented by the government to have a designated lane, and separated from other cars by concrete barriers \cite{Chua2026}. Throughout the years of implementation, a large number of commuters have utilized the system with an average of 183,000 daily passengers, climbing up to 320,000 during peak periods, with a ridership of nearly 67 million in 2025, up from 63.02 million in 2024 (Figure~\ref{fig:bg-ridership}) \cite{DOTr2025Ridership}. With that amount of ridership, the EDSA Bus Carousel can be safely assumed to be one of the backbones of Metro Manila commuting.

\begin{figure}[htbp]
    \centering
    \includegraphics[width=0.95\textwidth]{Figures/bg_fig1_ridership.pdf}
    \caption[EDSA Carousel ridership]{EDSA Carousel ridership. (a) Annual ridership
    grew 5.79\% from 63.02M in 2024 to 66.67M in 2025. (b) Daily ridership in 2025
    averaged 182{,}655 passengers, with a single-day peak of 321{,}186 in April 2025
    --- approximately 1.76$\times$ the daily mean \cite{DOTr2025Ridership}.}
    \label{fig:bg-ridership}
\end{figure}
%%

Figure~\ref{fig:bg-corridor-map} shows the EDSA Carousel Southbound route from Monumento to PITX, the corridor this study is grounded in, together with the other public transport modes (jeepney, MRT, LRT, tricycle, UV/FX) that intersect it at each major stop.

\begin{figure}[htbp]
    \centering
    \includegraphics[width=0.75\textwidth]{Figures/bg_fig3_edsa_corridor_map.pdf}
    \caption[EDSA Carousel corridor map]{EDSA Carousel Southbound route
    (Monumento to PITX) with intersecting public transport modes at each
    stop. Authors' illustration, adapted from the group's defense
    presentation.}
    \label{fig:bg-corridor-map}
\end{figure}

Despite the gradual increase in demand, the EDSA Bus Carousel still faces a lot of significant operational issues that affect the quality of bus routing. Firstly, the LTFRB relied on a tedious manual monitoring process by deploying personnel at each stop, creating decisions through empirical observation. \cite{EDSApolicy2023}. In relation to that, it is also clear that bus bunching has not been resolved; most of the time, their designated dispatch schedule is not followed \cite{Ollero2024EDSA}. There is an acknowledged lack of a structured monitoring framework for operations (e.g., dispatching, dwell time, bus bunching), in which adjustments are often done only after recurring issues surface \cite{Tiglao2025}. Moreover, mechanical failures or bus breakdowns are also prominent in Metro Manila, which contributes to a stressful environment, highlighting that poor vehicle condition and the constant threat of mid-route breakdowns are significant chronic stressors for Filipino bus drivers \cite{Chua2026}. Furthermore, environmental stochasticity severely destabilizes operations: empirical studies on Philippine expressways show that increasing rainfall intensity significantly reduces average traffic speed and free-flow capacity \cite{TSSP_Rain2018}. This rainfall-impact evidence is drawn from a 2018 study of the North Luzon Expressway rather than the EDSA Busway, and is used here only as contextual motivation that weather materially affects Philippine road-traffic operations; the weather-disturbance generator in this study (Section~3.2.6) does not adopt this study's specific speed-reduction percentages, and EDSA-specific travel-time behavior is independently calibrated through the GEH/RMSE procedure described in Section~3.2.3.

\begin{figure}[htbp]
    \centering
    \includegraphics[width=0.7\textwidth]{Figures/bg_fig2_rainfall_traffic.pdf}
    \caption[Rainfall impact on traffic flow]{Reduction in average speed and
    free-flow capacity relative to the clear-weather baseline across PAGASA
    rainfall categories. Capacity loss grows sharply at heavy intensities
    ($-17.4\%$), whereas mean speed degrades more moderately ($-7.4\%$), adapted from \cite{TSSP_Rain2018}.}
    \label{fig:bg-rainfall}
\end{figure}

As Figure~\ref{fig:bg-rainfall} shows, the reduction in
average speeds are about 5.34\% under light rain conditions, 6.3\% under moderate rain
conditions, and 7.4\% under heavy rain conditions. The average volume decreased by 2.92\%,
10.62\%, and 12.39\% while the capacities reduced by 3.67\%, 7.6\%, and 17.44\% under light,
medium, and heavy rain conditions, respectively \cite{TSSP_Rain2018}. Severe weather events often initiate localized flooding at key segments like Muñoz and Monumento Stations or cause severe degradation, forcing operational halts and emergency speed restrictions that further affect the busway’s performance  \cite{PhilstarTyphoon2024, PIA_Emergency2023, BusRepair2023}. Given these operational hurdles, the bus dispatching process of the EDSA Carousel can be modeled as a Vehicle Scheduling Problem (VSP).

Traditionally, the VSP has been addressed through static scheduling, which utilizes fixed timetables and headways derived from historical averages \cite{Ceder2007}. However, these rigid, rule-based methods often fail in stochastic environments as they cannot adapt to real-time disruptions, necessitating a shift toward more dynamic, autonomous scheduling frameworks that can respond to the unpredictability of the corridor \cite{Ollero2024EDSA}. This motivated a transition toward data-driven methodologies. Initial studies utilized traditional machine learning (ML) for demand forecasting and traffic prediction \cite{Wang2017}, but these models are primarily limited to predictive tasks and lack the active control mechanisms required for real-time decision-making in a closed-loop system. The subsequent wave introduced Single-Agent Reinforcement Learning (SARL), which enabled autonomous control, yet these frameworks frequently suffer from the `curse of dimensionality' as the joint state and action spaces grow exponentially with the number of vehicles, leading to centralized bottlenecks that hinder scalability in large transit networks \cite{Chen2024Survey}. To overcome these limitations, Multi-Agent Reinforcement Learning (MARL) has emerged as the state-of-the-art approach, offering a decentralized execution framework that enables scalable, collaborative control among multiple buses to mitigate complex issues like bus bunching \cite{Wang2020Holding}.

Although recent studies have demonstrated the potential of MARL-based dynamic scheduling to improve transit adaptability and overall reliability, a critical research gap remains. Existing MARL bus scheduling models are primarily evaluated in controlled simulation settings and are typically tested against isolated disruptions. No prior work has evaluated these models under the combined, simultaneous severe disruptions such as heavy rain, mechanical breakdowns, traffic variations, and sudden passenger surges that portray real-world corridors like EDSA. Therefore, this study aims to evaluate a MARL-based dynamic bus scheduling model under these non-ideal, stochastic traffic conditions to mitigate bus bunching, minimize passenger waiting times, and improve overall dispatch efficiency on the EDSA Carousel.

%%MARQUEZ END
\clearpage
\section{Review of Related Work}
This section surveys prior work on data-driven bus and traffic scheduling using a funnel structure. It starts from the broad transition from static scheduling toward adaptive and learning-based control, narrows to single-agent reinforcement learning and its known scalability limitations in large transit systems, narrows further to Multi-Agent Reinforcement Learning (MARL) and the assumptions commonly used in existing MARL bus-scheduling studies, and finally identifies the robustness gap addressed in this study.

\subsection{From Static Rules to Machine Learning}
\textbf{Static scheduling} refers to the offline determination of bus departure times, headways, fleet allocations, and stop assignments based on historical demand and travel-time averages, fixed in advance of operations and not adjusted in real time in response to observed disturbances \cite{Ceder2007}. Once deployed, such schedules cannot react to live disturbances such as demand surges, congestion, weather-induced delays, or vehicle breakdowns. The poor performance of fixed timetables under such conditions is well documented: even-headway and schedule-based control degrade rapidly once travel-time and demand variability exceed the narrow band assumed at design time, allowing small headway perturbations to amplify into bunching \cite{Daganzo2009}. Static schedules therefore remain mathematically inadequate for stochastic traffic environments, where the governing quantities are random variables rather than deterministic constants. Under the specific non-ideal conditions this study targets, the failure modes differ by control strategy. A fixed timetable has no feedback mechanism at all, so once bunching begins nothing in the schedule corrects it. A local reactive rule that holds a bus based only on the gap to the bus ahead can partially correct bunching under ordinary congestion, but has no way to respond to a breakdown, since it observes only the forward gap and not the enlarged gap a failed bus leaves behind it. A more globally aware reactive rule that accounts for both the forward and backward gap improves on this, but still follows a fixed, pre-specified rule rather than a learned response, so it cannot adapt its behavior to the heavier-tailed delays that severe weather introduces. These limitations motivated the transition toward more adaptive and data-driven scheduling methodologies.

Machine Learning (ML) techniques subsequently emerged as important tools for improving operational responsiveness in transportation systems through demand estimation, traffic forecasting, predictive dispatching, and adaptive holding strategies \cite{huang2019,patil2025}. Unlike static scheduling approaches, ML-assisted frameworks incorporate continuously updated operational information to improve short-term decision support under uncertain traffic conditions \cite{huang2019}. Recent applications to transit operations have demonstrated measurable operational gains. Barrera Hernandez et al. applied passenger-demand forecasting to support a heuristic dispatcher for bus operations in Bogotá, reporting reductions in passenger waiting time and operational cost relative to fixed timetables \cite{Barrera2025Optimization}. Similarly, Usman et al. utilized ML-based density estimation to improve real-time responsiveness in public transit systems \cite{Usman2025ML}. These studies demonstrate the growing role of ML-assisted methodologies in adaptive transit scheduling and operational planning.

Despite these advances, many ML-assisted transit-control frameworks still rely heavily on manually engineered intervention logic and externally designed dispatch heuristics \cite{alexandre2023}. In several predictive-control and heuristic-dispatch systems, forecasting components improve operational visibility, but the resulting operational responses are still governed by manually specified holding rules, threshold conditions, or predefined intervention strategies \cite{pang2019,verbich2021}. In the framework proposed by Barrera Hernandez et al., for example, the forecasting model informs the dispatcher but does not autonomously determine operational actions such as holding duration, dispatch timing, or stop-skipping behavior \cite{Barrera2025Optimization}. Consequently, while ML-assisted scheduling frameworks improve anticipatory planning and adaptive responsiveness, many existing implementations still stop short of establishing a fully integrated closed-loop control framework capable of autonomously optimizing sequential dispatch decisions over time.

\begin{table}[htbp]
\centering
\caption{Progression from static rules to reinforcement learning along the dimensions that determine real-time control capability.}
\label{tab:control_paradigms}

\footnotesize
\renewcommand{\arraystretch}{1.3}

\begin{tabularx}{\textwidth}{
>{\RaggedRight\arraybackslash}p{2.8cm}
>{\RaggedRight\arraybackslash}p{3.5cm}
>{\RaggedRight\arraybackslash}p{3.5cm}
>{\RaggedRight\arraybackslash}X
}
\toprule
\textbf{Paradigm} & 
\textbf{Adapts to real-time conditions?} & 
\textbf{Produces a dispatching action?} & 
\textbf{How the action is determined} \\
\midrule

Static / heuristic scheduling & 
No, fixed offline from historical averages & 
Yes, but actions are predetermined before operation & 
Pre-set timetables or manually designed heuristic rules that do not adapt to live corridor conditions \cite{Ceder2007,Daganzo2009} \\

Machine-learning assisted scheduling & 
Yes, forecasts update continuously using observed operational data & 
Partially; predictions support decisions, but actions are still governed externally & 
Forecasting models estimate demand or traffic conditions while separate heuristic dispatch rules determine the operational response \cite{Barrera2025Optimization,Usman2025ML,pang2019} \\

Reinforcement learning & 
Yes, policies continuously condition on live observed system states & 
Yes, the dispatching action is produced directly by the learned policy & 
The control policy learns through interaction with the environment to maximize long-term operational performance \cite{Momenikorbekandi2023Intelligent,Rodriguez2023Cooperative,Wangsun} \\

\bottomrule
\end{tabularx}
\end{table}

The decisive distinction is the final column: only reinforcement learning makes the dispatching action itself a learned output. This study instantiates the reinforcement-learning row as a parameter-sharing Double Deep Q-Network under centralized training with decentralized execution (Chapter~3).
\FloatBarrier
\subsection{Single-Agent Reinforcement Learning and Its Limitations}
Reinforcement Learning (RL) closes the prediction--control gap by learning a \textit{policy}: a mapping from observed system states to actions, optimized through trial-and-error interaction with the environment to maximize a long-horizon reward signal \cite{Momenikorbekandi2023Intelligent}. Unlike a forecaster, an RL agent does not just predict what will happen --- it decides what to do. The reward signal closes the loop: the agent's action changes the environment, the environment returns a new state and a scalar reward, and the policy is updated so that better actions become more likely the next time a similar state is encountered.

When RL is applied to a multi-vehicle problem such as bus-fleet control, the simplest formulation treats the whole fleet as one entity controlled by a single agent. This is \textbf{Single-Agent Reinforcement Learning (SARL)}: one centralized policy observes the full system state and emits a joint action covering every bus at once. SARL has been applied directly to bus holding control, including the STDH-DQN approach of Zhao et al.~\cite{Zhao2022STDH}, which uses a self-attention state encoder over spatial-temporal AVL features, and the more recent SA-DRL approach of Zhang and Zheng~\cite{Zhang2025SADRL}, which encodes categorical identifiers (vehicle ID, station ID, trip ID) as state features and reports favorable results against MARL baselines on a bidirectional timetabled line. SARL is therefore a real and competitive option in this space; the question is not whether it works but where it begins to strain.

SARL faces three scalability limitations in multi-bus control, each documented in the bus-scheduling and traffic-control literature. These limitations motivate the move to a multi-agent formulation in the next subsection; each has a corresponding design response that MARL provides.

\begin{enumerate}
\item \textbf{State-space dimensionality growth.} A SARL formulation of $N$ buses on $M$ stops requires a global state $s \in \mathbb{R}^{N \cdot d}$, where $d$ is the per-bus feature count (position, headway, load). For the EDSA Carousel sub-corridor studied here, with $M = 24$ stops and a fleet of $N \approx 12$--$30$ buses, $\dim(s)$ reaches roughly $96$--$240$. The function approximator must also learn approximate permutation invariance over buses, the policy should not depend on which arbitrary index a bus is assigned, a property not enforced by default in neural networks \cite{Wang2023MultiObj}.

\item \textbf{Joint action-space combinatorial growth.} A SARL controller emits a joint action $a = (a_1, \ldots, a_N) \in A^N$. With the hybrid action space adopted in this study ($|A_i| = 5 \times 2 = 10$, combining a 5-bin holding-strength discretization with a binary skip decision), a SARL controller managing $N = 20$ buses operates over a joint action space of $10^{20}$ combinations. Although deep RL policies do not enumerate these combinations directly --- they sample from a parameterized distribution over actions --- the credit-assignment burden remains: a single scalar reward must be attributed across all $N$ simultaneous component actions, and this attribution problem worsens as $N$ grows. MARL with shared policies factorizes the action selection into $N$ independent decisions of size $10$, each with its own attributable reward signal.

\item \textbf{Sample inefficiency under non-stationarity.} SARL treats the fleet as a single dynamical system; every demand change, breakdown, or weather event triggers re-adaptation of the entire policy. MARL with parameter sharing benefits from \textit{experience aggregation}: every action by every bus generates one training sample for the same shared network, so each bus's experience accelerates learning for all of them \cite{Christianos2021PS,Yu2022MAPPO}.
\end{enumerate}

Recent SARL-for-bus-control work that competes successfully against MARL baselines does so partly by hand-engineering categorical agent-identity features into the state representation \cite{Zhang2025SADRL}. This is informative because it suggests that agent identity remains an important representational issue in SARL formulations even at the high end of reported performance --- the dimensionality concern in item~(1) above is partially addressed by smuggling identity through input features rather than by a change in policy class. The fix is workable on a fixed corridor but does not transfer cleanly when fleet size or route topology changes.

We note a counter-position. Zhang and Zheng~\cite{Zhang2025SADRL} report SA-DRL matching or exceeding MARL baselines on bidirectional timetabled lines, attributing this to data-imbalance issues in MARL when some agents see few episodes. The EDSA Carousel sub-corridor studied here is operated as a high-frequency BRT service along a defined directional segment, closer to the single-direction loop-line assumptions where MARL has been shown to excel, which motivates the MARL choice here while acknowledging this caveat.

To situate the MARL literature reviewed in the next subsection within the broader ML and SARL landscape, Table~\ref{tab:ml_sarl_coverage} extends the paradigm comparison of Table~\ref{tab:control_paradigms} with a disturbance-coverage column, using the same D/S/T/W/B notation as Table~\ref{tab:marl_performance}.

\begin{table}[htbp]
\centering
\caption{Disturbance coverage across ML and SARL vehicle-scheduling studies, preceding the MARL-specific comparison in Table~\ref{tab:marl_performance}.}
\label{tab:ml_sarl_coverage}

\footnotesize
\renewcommand{\arraystretch}{1.3}

\begin{tabularx}{\textwidth}{
>{\RaggedRight\arraybackslash}p{2.6cm}
>{\RaggedRight\arraybackslash}p{2.0cm}
>{\RaggedRight\arraybackslash}X
>{\Centering\arraybackslash}p{1.6cm}
}

\toprule
\textbf{Paper} &
\textbf{Paradigm} &
\textbf{Method} &
\textbf{Disturbances} \\
\midrule

Wang et al.~\cite{Wang2017} &
ML (data-driven) &
Bus scheduling model incorporating time-dependent traffic and demand &
D \\

Barrera Hernandez et al.~\cite{Barrera2025Optimization} &
ML-assisted (heuristic dispatcher) &
Passenger-demand forecasting supporting a heuristic dispatcher &
D \\

Zhao et al.~\cite{Zhao2022STDH} &
SARL &
STDH-DQN; self-attention state encoder over spatial-temporal AVL features &
D, T \\

Zhang and Zheng~\cite{Zhang2025SADRL} &
SARL &
SA-DRL; categorical identity features (vehicle, station, trip ID) &
D, T \\

Verbich and El-Geneidy~\cite{verbich2021} &
Heuristic (non-MARL) &
Dynamic transit control under severe weather and vehicle breakdowns &
W, B \\

\bottomrule
\addlinespace[2pt]
\multicolumn{4}{@{}p{\textwidth}@{}}{\footnotesize D~=~stochastic demand, S~=~demand surges, T~=~traffic-speed perturbation, W~=~weather-induced heavy-tail delay, B~=~discrete vehicle breakdown. Disturbance coverage for the ML/SARL rows follows the characterization given by the RTC panel during the proposal review; unlike Table~\ref{tab:marl_performance}, these classifications have not been independently re-verified against each paper's full text.} \\
\end{tabularx}
\end{table}

The funnel is now complete: no ML or SARL study covers W or B, and among MARL studies (Table~\ref{tab:marl_performance}), only Verbich and El-Geneidy's heuristic controller~\cite{verbich2021} addresses both --- and it is explicitly non-MARL. Patil et al.~\cite{Patil2025Conformal} similarly validate weather-induced travel-time distributions but do not address bus control at all; their contribution to this study is the lognormal parameterization used by the weather-disturbance generator (Section~3.2.6), not a bus-control baseline. No prior study, ML, SARL, or MARL, combines W and B coverage with an actual MARL bus-scheduling controller, which is the specific gap this study fills.

\subsection{Multi-Agent Reinforcement Learning}

\begin{figure}[htbp]
    \centering
    \includegraphics[clip, trim=0.1cm 0.7cm 0.1cm 0.1cm,width=\textwidth]{Figures/rrl_fig1_sarl_vs_marl (3).pdf}
    \caption[SARL vs MARL architecture]{Comparison of single-agent and
    multi-agent formulations. (a) SARL: a centralized
    policy $\pi_\theta(s)$ consumes the concatenated global state
    $s \in \mathbb{R}^{N \cdot d}$ and outputs a joint action covering every bus.
    (b) MARL with parameter sharing: a single set of weights $\theta$ is shared
    across $N$ agents, each acting on its own local observation $o_i$. Authors' illustration.}
    \label{fig:sarl-vs-marl}
\end{figure}

In both panels of Figure~\ref{fig:sarl-vs-marl}, the per-bus state (denoted $s_{i,t}$ in the formal MDP notation of Section~\ref{subsec:state-space}, and shown in the figure as the local observation $o_i$) encodes the bus's current position, forward and backward headways, onboard load, and queue length at its current stop, as defined in full in Section~\ref{subsec:state-space}. The action $a_{i,t}$ is the holding-strength and stop-skipping decision the controller emits for that bus, defined in Section~\ref{subsec:action-space}. In the SARL panel (a), a single centralized network ingests all $N$ per-bus state vectors concatenated into one global state $s \in \mathbb{R}^{N \cdot d}$ and outputs all $N$ actions simultaneously; in the MARL panel (b), the same shared network weights $\theta$ instead process each bus's local state independently, so each agent acts on only its own observation rather than the concatenated global one.

Multi-Agent Reinforcement Learning (MARL) addresses the three limitations above by decomposing decision-making across multiple agents that share the environment. Each limitation has a direct design counterpart in MARL: state-space growth is contained by giving each agent only its own \textit{local} observation rather than the global state; action-space growth is factorized by having each agent select only its own action; and the sample-inefficiency problem is mitigated by sharing one network's parameters across all agents, so that every agent's experience updates the same policy. The problem is typically formalized as a Markov Game built on an underlying Markov Decision Process, with each agent maintaining its own local observation and acting under a shared or coordinated policy \cite{Li2023Departure,Sun2024Graph,Zuo2026AGV,Huang2025Joint,Bokade2023MARL}.

This decomposition introduces a new difficulty: \textit{partial observability}. No individual bus can see the full corridor state; it knows only its own headway, load, and immediate surroundings. A fully decentralized scheme in which each agent both trains and acts on its local observation alone would be unstable, because the other agents' simultaneously evolving policies make the environment non-stationary from any one agent's viewpoint --- the same local state can transition to different outcomes depending on what the other agents happen to be learning at that moment.

To address this, most modern formulations use \textbf{Centralized Training with Decentralized Execution (CTDE)}, illustrated in Figure~\ref{fig:ctde}. During training, agents are allowed to access system-wide information and optimize a shared objective; this stabilizes learning despite the non-stationarity introduced by simultaneously adapting agents. At deployment, each bus acts only on its own local observation, which is what makes the framework practically deployable on a real fleet without requiring a centralized real-time controller \cite{Cao2022Train,Wang2024MultiAGV,Nie2025CMRM,Xu2026Hierarchical}.



\begin{figure}[htbp]
    \centering
    \includegraphics[width=\textwidth]{Figures/rrl_fig2_ctde (2).pdf}
    \caption[Centralized training with decentralized execution]{(a) During training, all agents share one reward
    signal and one set of policy weights $\theta$, which are updated centrally
    using joint-state information. (b) At deployment, the weights are frozen and
    each bus runs an identical copy of $\pi_\theta$ on its own local observation
    $o_i$ alone, so no centralized real-time controller is required. Authors' illustration.}
    \label{fig:ctde}
\end{figure}

In the bus-scheduling setting, a parameter-sharing CTDE architecture is fleet-size invariant: the same trained policy applies regardless of how many buses are currently in service, because every bus is an instance of the same agent class acting on its own local headway, occupancy, and queue observations. Adding or removing a bus does not invalidate the trained model. Christianos et al.~\cite{Christianos2021PS} demonstrate that parameter sharing scales effectively to large agent populations and provides empirical support for the fleet-size invariance claim.

Recent progress has combined MARL with deep learning algorithms --- Deep Q-Networks (DQN), Proximal Policy Optimization (PPO) and its multi-agent variant MAPPO, and Graph Attention Networks (GAT) --- to handle high-dimensional observation spaces and complex interactions. These have been applied across vehicle scheduling, signal control, and logistics, generally reporting improvements in routing efficiency, congestion, and throughput in simulation \cite{Li2023Departure,Cao2022Train,Wang2024MultiAGV,Sun2024Graph,Zuo2026AGV,Huang2025Joint,Bokade2023MARL}.

\subsection{MARL Applied to Bus Scheduling}
\label{subsec:marl-applied}
Within bus scheduling specifically, the MARL literature has accumulated along a recognizable trajectory: first establishing that holding control could be learned at all, then enriching the action and objective space, and most recently beginning to examine policy behavior under disturbance. This trajectory is what the present study extends.

The foundational result is due to Wang and Sun~\cite{Wang2020Holding}, who showed that a cooperative MARL framework could learn an effective bus-holding policy on a single-line corridor and outperform classical headway-equalization rules under idealized stochastic demand. Chen et al.~\cite{Chen2016MARLBus} provided an earlier proof-of-concept in the same vein on a corridor with multi-agent reinforcement learning over local headway observations. These early results established the viability of the approach but were limited to the holding action alone, restricting the policy's leverage once bunching had already developed.

Subsequent work enriched the formulation along two axes. Rodriguez et al.~\cite{Rodriguez2023Cooperative} extended the action space by combining holding with stop-skipping in a cooperative DDQN framework on a simulated transit corridor with Poisson-distributed passenger arrivals and Gaussian inter-stop travel times, reporting reductions in average passenger waiting time relative to an even-headway holding baseline. Wang and Sun~\cite{Wang2023MultiObj} extended the objective space, scaling MARL to multi-line services with a shared corridor and jointly optimizing headway regularity against operational cost under controlled simulated demand and travel-time profiles.

Recent work has increasingly examined policy behavior under disturbance. Wang and Sun's IQNC-M~\cite{Wangsun} introduced implicit quantile networks and meta-learning for risk-sensitive bus holding, evaluating performance under demand surges and traffic-speed perturbations but not under weather-induced heavy-tailed delays or discrete vehicle breakdowns. In adjacent settings, Katzilieris et al.~\cite{Katzilieris2026MARL} integrated MADDPG-based traffic signal control with dynamic bus lane access management on a simulated arterial network with deterministic signal cycles; Liu et al.~\cite{Liu2023DRLHolding} reported DRL-based holding control under stochastic travel time and demand; and Shi et al.~\cite{Shi2022DistDRL} reported a distributed DRL bus control system in a connected environment with breakdown events modeled but without weather effects.

Across this trajectory, the simulation environments share three idealizing assumptions: passenger arrivals are stationary Poisson processes, inter-stop travel times are symmetric and well-behaved (Gaussian or similar), and the fleet experiences no discrete failure events. Reported performance gains are therefore conditional on those assumptions holding at deployment. The EDSA Carousel sub-corridor studied here violates each of them. Travel times are highly non-stationary, with tail behavior driven by tropical-cyclone and heavy-rainfall events that suspend operations entirely on the most severe days \cite{DOTr2020Suspension}; passenger demand swings sharply across peak intervals, with single-day ridership peaks roughly $1.8\times$ the daily average \cite{DOTr2025Ridership}; and tight headways mean that any individual vehicle breakdown propagates through the corridor as a downstream demand-absorption problem rather than as an isolated event. A controller whose reported gains rest on the idealizing assumptions above has no validated behavior on a corridor that violates all three at once.

\begin{table}[htbp]
\centering
\caption{Disturbance coverage across MARL bus-scheduling literature.}
\label{tab:marl_performance}

\footnotesize
\renewcommand{\arraystretch}{1.3}

\begin{tabularx}{\textwidth}{
>{\RaggedRight\arraybackslash}p{2.5cm}
>{\RaggedRight\arraybackslash}p{3.0cm}
>{\RaggedRight\arraybackslash}p{2.8cm}
>{\Centering\arraybackslash}p{2.2cm}
>{\RaggedRight\arraybackslash}X
}

\toprule
\textbf{Paper} &
\textbf{Algorithm} &
\textbf{Environment} &
\textbf{Disturbances} &
\textbf{Reported gain} \\
\midrule

Wang \& Sun (2020)~\cite{Wang2020Holding} &
Cooperative MARL holding &
Single-line, idealized &
D &
Bunching reduction vs.\ headway rule \\

Chen et al.\ (2016)~\cite{Chen2016MARLBus} &
Early MARL holding &
Corridor, idealized &
D &
Proof-of-concept gains \\

Rodriguez et al.\ (2023)~\cite{Rodriguez2023Cooperative} &
Cooperative DDQN (hybrid action) &
Chicago, mesoscopic &
T &
Wait time vs.\ headway-control baseline \\

Wang \& Sun (2023)~\cite{Wang2023MultiObj} &
Multi-objective MARL &
Shared corridor, idealized &
D &
Multi-line bunching reduction \\

Wang \& Sun (2023)~\cite{Wangsun} &
IQNC-M (distributional) &
Real-world bus services &
S, T &
Stable control under extreme events \\

Katzilieris et al.\ (2026)~\cite{Katzilieris2026MARL} &
MADDPG &
Athens arterial, SUMO &
T &
Throughput improvement \\

Liu et al.\ (2023)~\cite{Liu2023DRLHolding} &
DRL holding &
Synthetic single line &
D, T &
Wait time vs.\ no-control \\

Shi et al.\ (2022)~\cite{Shi2022DistDRL} &
Distributed DRL &
Connected environment &
D, B &
Integrated control gains \\

\textbf{This study} &
\textbf{DDQN + PS (CTDE)} &
\textbf{Calibrated EDSA SUMO} &
\textbf{S, T, W, B} &
\textbf{To be evaluated} \\

\bottomrule
\addlinespace[2pt]
\multicolumn{5}{@{}p{\textwidth}@{}}{\footnotesize D~=~stochastic demand, S~=~demand surges, T~=~traffic-speed perturbation, W~=~weather-induced heavy-tail delay, B~=~discrete vehicle breakdown.} \\
\end{tabularx}
\end{table}

\clearpage
Only Shi et al.~\cite{Shi2022DistDRL} carries a B (breakdown) entry in Table~\ref{tab:marl_performance}. Cao et al.~\cite{Cao2022Train}, which also models discrete vehicle failures, is deliberately excluded from this count: their MARL application is to \textit{train} rescheduling, not bus scheduling, so it does not belong in a table scoped to MARL bus-control literature. Verbich and El-Geneidy~\cite{verbich2021} likewise model breakdowns but use heuristic, non-MARL control (Table~\ref{tab:ml_sarl_coverage}), so they are excluded for the same reason. Among MARL bus-scheduling studies specifically, Shi et al. remains the only one to model discrete breakdowns.

Table~\ref{tab:marl_performance} summarizes what each study evaluated, what disturbances it modeled, and what it reported. No prior MARL bus paper has evaluated controller performance under a framework that incorporates all four disturbance classes documented in the EDSA corridor's operational profile such as demand surges, traffic-speed variability, heavy-tailed weather-induced travel-time delays, and discrete breakdowns in a calibrated real-corridor microsimulation. Of these four, weather-induced heavy-tailed delays and discrete breakdowns are the two that no prior MARL bus-scheduling study has modeled at all; their combination with demand surges and a calibrated traffic-speed baseline is the specific disturbance profile introduced here. The algorithmic choice this study adopts to address that combined-disturbance setting, DDQN with parameter sharing under CTDE is justified in detail in Chapter~3 (Section~\ref{subsec:proposed-algorithm}), alongside the rationale for not using actor-critic or distributional alternatives. The action space the agents operate on is correspondingly small and discrete: at each control event, an agent selects a holding strength from a five-bin discretization $\alpha \in \{0.0, 0.1, 0.2, 0.3, 0.4\}$ (multiplied by a maximum holding duration to yield seconds held) combined with a binary decision to skip the next stop, giving $|A_i| = 10$ actions per event. The full action-space rationale, including why discretization is preferred over a continuous holding parameter, is given in Chapter~3 (Section~\ref{subsec:action-space}).



Under the non-ideal evaluation regime the weather generator (W) \textbf{replaces} the traffic-speed generator (T) as the source of inter-stop travel-time variability rather than acting alongside it. Both are anchored to the same per-segment empirical mean $\mu$, and the weather coefficient of variation ($\eta \ge 0.3$) already spans and exceeds the $\pm 20\%$ range produced by traffic-speed scaling (Section~3.2.6). T therefore governs travel-time stochasticity under the ideal regime and W governs it under the non-ideal regime; the two are not composed on the same segment.



\subsection{The Disturbance Gap}
\label{subsec:disturbance-gap}
The idealized assumptions in Table~\ref{tab:marl_performance} limit how confidently the reported MARL gains can be expected to transfer to real urban corridors. Real corridors violate all three assumptions: demand exhibits non-stationary surges, weather-induced travel times are skewed and heavy-tailed, and individual buses break down stochastically. Wang and Sun themselves identify this as central, noting that existing MARL bus-control studies overlook performance under perturbations and that this becomes critical when transferring models toward real-world deployment \cite{Wangsun}. Their IQNC-M framework is among the first to address demand surges and traffic-speed perturbations in bus holding control, but weather-induced heavy-tailed delays and discrete vehicle breakdowns remain outside its tested scenario set.

Two considerations explain why the gap has persisted:

\begin{enumerate}
\item \textbf{Training instability under heavy-tailed returns.} Prior MARL bus-scheduling studies have typically used Poisson arrivals and Gaussian travel-time noise because these distributions generally produce more statistically stable training behavior than heavy-tailed disturbances. In reinforcement learning literature, heavy-tailed returns and stochastic gradients are known to introduce convergence difficulties and unstable optimization behavior during policy learning \cite{Henderson2018Matters,Agarwal2021Precipice}.
\clearpage
\item \textbf{Lack of temporally aligned anomaly data.} Historical severe-weather records are rarely aligned with operational bus data at the per-segment granularity needed for direct parameter estimation, making it difficult to fit an empirically grounded heavy-tailed disturbance distribution from corridor-specific data alone, especially given the short operational observation windows typically available. \textit{This study addresses the obstacle by adopting the lognormal travel-time parameterization that Patil et al.\ \cite{Patil2025Conformal} validated against INRIX freeway data via the Kolmogorov-Smirnov test, and then sweeping the disturbance intensity across a range of validated and extrapolated CV values rather than fitting a single point estimate. This characterizes controller behavior as a function of disturbance magnitude and avoids reliance on a corridor-specific severe-weather sample whose variance cannot be reliably resolved from a short observation window.}
\end{enumerate}

\subsection{Significance of Closing the Gap}
This study contributes the first characterization of MARL bus-scheduling performance across a swept range of disturbance intensities on a calibrated real-corridor simulation, with the disturbance set covering heavy-tailed weather-induced travel-time delays and discrete bus breakdowns that no prior MARL bus-scheduling study has modeled. Understanding policy performance under disturbance matters for two reasons. First, deploying MARL bus-scheduling models on the basis of ideal-simulator benchmarks alone leaves their real-world failure modes unknown until deployment, which is an unacceptable risk profile for a system that carries hundreds of thousands of daily passengers. Second, the relevant performance question for transit operators is not only average-case waiting time under quiet conditions but how a control policy degrades as a disruption intensifies, a characterization that prior MARL bus-scheduling work does not report. Establishing a reproducible stress-testing protocol pairing a calibrated corridor simulation with configurable demand, weather, and breakdown generators therefore contributes both an empirical characterization of performance under disturbance and a reusable evaluation framework for subsequent MARL work in this domain \cite{Li2023Departure,Wang2024MultiAGV,Che2024Recharging,Huang2025Joint,Bokade2023MARL}.")

% Line Numbering on the left side
\usepackage[left]{lineno}
% ------------------------------------------------------------------------------------------


% ------------------------------- TITLE ---------------------------------------

% MODIFY THIS SECTION

\title{An Evaluation of Multi-Agent Reinforcement Learning for Dynamic Bus Scheduling Under Non-Ideal Conditions}

% Update names of proponents and supervisor in title.tex

% --------------------------------------------------------------------------------


% ----------------------------------------------------------------------------------
% THIS IS THE BEGINNING OF THE MANUSCRIPT PROPER
% ----------------------------------------------------------------------------------
%TC:ignore
\begin{document}
%TC:ignore
\onehalfspacing
\begin{titlepage}

\newcommand{\HRule}{\rule{\linewidth}{0.5mm}} % Defines a new command for the horizontal lines, change thickness here

%----------------------------------------------------------------------------------------
%	LOGO SECTION
%----------------------------------------------------------------------------------------

% DO NOT MODIFY THIS SECTION!!!

\center % Center everything on the page


\includegraphics[width=3.5cm]{Figures/UST_logo.jpg}

\includegraphics[width=2.8cm]{Figures/ustenglogo.png}\\[2cm] % Include a department/university logo - this will require the graphicx package
 
%----------------------------------------------------------------------------------------

%----------------------------------------------------------------------------------------
%	TITLE SECTION
%----------------------------------------------------------------------------------------
\makeatletter
\HRule \\[0.4cm]
{ \huge \bfseries \@title}\\[0.4cm] % Title of your document
\HRule \\[1.7cm]

%----------------------------------------------------------------------------------------
%	HEADING SECTIONS
%----------------------------------------------------------------------------------------

% DO NOT MODIFY THIS SECTION!!!

\text{A thesis submitted in partial fulfilment}\\[0.2cm]
\text{of the requirements for the degree of}\\[0.2cm]
\text{Bachelor of Science in Electronics Engineering}\\[1.5cm] 
\text{Department of Electronics Engineering}\\[0.2cm] 
\text{Faculty of Engineering}\\[0.2cm] 
\text{University of Santo Tomas}\\[1.7cm] 

%----------------------------------------------------------------------------------------
%	AUTHOR SECTION
%----------------------------------------------------------------------------------------

% MODIFY THIS SECTION: Names of proponents and supervisor

\begin{minipage}{0.4\textwidth}
\begin{flushleft} %\large
\emph{Proponents:} \\  %Enter the name of the authors here
Badal, Khalil T. \\[0.2em]
Lopez, Elijah David B.\\[0.2em]
Mananguit, Jared Asher N. \\[0.2em]
Marquez, Rence Joseph L. \\[0.2em]
Medenilla, Jose Anton M.\\[0.2em]
\end{flushleft}
\end{minipage}
~
\begin{minipage}{0.4\textwidth}
\begin{flushright} %\large
\emph{Research Supervisor:}\\
Asst. Prof. Kanny Krizzy D. Serrano, MSc \\[0.2em] % Edit the name of your research supervisor
\end{flushright}
\end{minipage}\\[1.5cm]
\makeatother

%----------------------------------------------------------------------------------------
%	DATE SECTION
%----------------------------------------------------------------------------------------
\begin{minipage}{0.4\textwidth}
\begin{flushleft} %\large
\emph{Date: }{ \today} % Date: change the \today to a set date if you want to be precise
\end{flushleft}
\end{minipage}
~
\begin{minipage}{0.4\textwidth}
\begin{flushright}%\large
\emph{Word Count: }{\automatedwordcount} % Word count of your manuscript starting from the introduction section onwards
% NOTE: You may use the 'Word Count' feature of Overleaf found in the Menu on the upper left corner of the window 
\end{flushright}
\end{minipage}\\[1cm]
%\makeatother

\vfill % Fill the rest of the page with whitespace

\end{titlepage} % Putting the title page

% ------------------------------- ABSTRACT ---------------------------------

% MODIFY THIS SECTION
\begin{comment}
\section*{Abstract} \large 
The abstract provides a concise summary of the project, highlighting its key objectives, methodologies, findings, and implications. It serves as a brief yet comprehensive overview, enabling readers to quickly grasp the essence of the work and its significance. (1) Include a brief introduction outlining the purpose and scope of the project. (2)Summarize the main objectives or research questions addressed in the project. (3) Describe the methodologies or approaches employed to investigate the objectives or address the research questions. (4) Highlight the key findings, results, or outcomes obtained through the research or project activities. (5) Discuss the implications, significance, or potential applications of the findings or outcomes in relevant fields or contexts. (6) Conclude with a statement emphasizing the contribution of the project to the existing body of knowledge and its relevance to stakeholders or interested parties.

[Optional: Include keywords or phrases relevant to the subject matter of the document/project for indexing and searchability purposes.]

[Optional: Provide acknowledgments or credits to individuals, organizations, or funding sources that supported the project.]

[Optional: Include a statement indicating any limitations or areas for further research identified during the course of the project.]

[Optional: Offer recommendations or suggestions for future initiatives or studies based on the findings or outcomes of the project.]

[End with any final remarks or closing thoughts that reinforce the significance of the project and its potential impact.] 

\clearpage
% --------------------------------------------------------------------------------
\end{comment}
% -------------------- ACKNOWLEDGEMENTS ------------------------

% MODIFY THIS SECTION
\begin{comment}
\section*{Acknowledgement} \large

% ------------------- You can remove this part --------------------------
\textbf{\textit{[NOTE]: This is just a sample acknowledgement statement. Kindly modify this accordingly.}}
% -----------------------------------------------------------------------

This work could not have been accomplished successfully without the help of many people.

We would like to express our sincere gratitude to all those who have contributed to the completion of this thesis. Without their support, encouragement, and assistance, this endeavor would not have been possible.

First and foremost, we are deeply indebted to our supervisor/advisor [Name], whose guidance, expertise, and unwavering support have been invaluable throughout this journey. Their insightful feedback, constructive criticism, and encouragement have significantly shaped the development of this thesis.

We extend our heartfelt appreciation to the members of our thesis committee, [Name(s)], for their valuable input, suggestions, and scholarly insights. Their expertise and dedication have enriched the quality of this work and broadened our perspective on the subject matter.

We are grateful to [Name(s) of Academic Department/Institution] for providing access to resources, facilities, and research opportunities essential for the completion of this thesis. Their institutional support has been instrumental in facilitating the research process and academic endeavors.

We would like to thank our colleagues and peers for their camaraderie, encouragement, and intellectual exchange throughout this academic endeavor. Their friendship, support, and collaborative spirit have made this journey both rewarding and memorable.

Special thanks are due to our family for their unwavering love, encouragement, and understanding during moments of academic challenges and triumphs. Their unconditional support and belief in our abilities have been a constant source of motivation and inspiration.

Last but not least, we express our profound gratitude to all individuals who participated in this research study or provided assistance in any form. Their contributions have enriched the findings and added depth to the scholarly discourse explored in this thesis.

In conclusion, we are deeply grateful to everyone who has played a part, big or small, in the completion of this thesis. Your support, guidance, and encouragement have been instrumental in this academic achievement, and for that, I am truly thankful.

\end{comment}

\section*{AI Use Disclosure} \large 


The authors acknowledge the use of artificial intelligence (AI)-assisted tools during the preparation of this manuscript. ChatGPT and Grammarly were used primarily as editorial, organizational, and formatting support tools during the writing and manuscript-preparation process.

Grammarly was used throughout all chapters for grammar checking, spelling correction, punctuation review, and sentence-level clarity suggestions. ChatGPT was used selectively for proofreading and language refinement of draft prose to improve readability, reduce informal phrasing, and smooth transitions between sections. It was also used to provide organizational feedback on manuscript structure, particularly in refining the funnel progression of the Review of Related Work, summarizing disturbance coverage across studies already collected by the authors, checking consistency of terminology, notation, section numbering, and symbol usage across Chapters~1--3, and identifying undefined terms, missing bibliography entries, and citation inconsistencies prior to compilation.

In addition, AI-assisted tools were used for document-preparation and Latex-related support tasks, including table formatting, float placement adjustments, bibliography troubleshooting, citation-key checking, and resolving compilation issues encountered during Overleaf manuscript assembly.

AI tools were employed strictly as editorial, organizational, and formatting assistants during the preparation of this manuscript. They were not used to generate research findings, conduct simulations, or produce technical conclusions. All core research decisions,including the proposed simulation framework, methodological direction, disturbance-modeling approach, evaluation design, and interpretation of the reviewed literature, were determined and validated by the authors. All AI-generated suggestions were reviewed, revised where necessary, and approved prior to inclusion. The authors assume full responsibility for the accuracy, originality, and integrity of the proposed work.





% --------------------------------------------------------------------------------


% -------------------- OTHER SECTIONS -------------------------------

% DO NOT MODIFY THIS SECTION

% Just modify the contents of the individual sections in their
% respective .tex files

\tableofcontents % Inserts the table of contents
\listoffigures % Inserts the list of figures
\listoftables % Inserts the list of tables


% --------------------------------------------------------------------------------
% Start line numbering
\linenumbers

% NOTE: Remove this line for the final, revised manuscript

% --------------------------------------------------------------------------------
%TC:endignore
\chapter{Introduction and Literature Review}

% =============================================================
% PART A — INTRODUCTION
% =============================================================

%% MARQUEZ START
\section{Background of the Study}
The public transportation system has long been a commodity to people; it operates on different terrains, has different forms, and has predetermined routes depending on their geographical location and the needs of commuters \cite{Tironi2018}. With an increase in urbanization over time, the demand for transportation has also increased. Despite the predicted decline in urbanization rates in Southeast Asia, the Philippines is experiencing the opposite \cite{UN_WUP2018}. Due to the increased urbanization and the continuous increase in population (especially in urban areas), people start looking for different alternatives in transportation \cite{Cervero2000}. With the increasing number of motor vehicles on the road, the traffic congestion worsens, which leads to commuters finding other means to quickly reach their destination. In Metro Manila, severe traffic congestion, stemming from insufficient road infrastructure and poor traffic mitigation policies, has resulted in a highly inefficient public transport system \cite{Spatio2026}. As a response, the Philippine government launched the EDSA Bus Carousel in 2020. It is the only busway system in Metro Manila implemented by the government to have a designated lane, and separated from other cars by concrete barriers \cite{Chua2026}. Throughout the years of implementation, a large number of commuters have utilized the system with an average of 183,000 daily passengers, climbing up to 320,000 during peak periods, with a ridership of nearly 67 million in 2025, up from 63.02 million in 2024 (Figure~\ref{fig:bg-ridership}) \cite{DOTr2025Ridership}. With that amount of ridership, the EDSA Bus Carousel can be safely assumed to be one of the backbones of Metro Manila commuting.

\begin{figure}[htbp]
    \centering
    \includegraphics[width=0.95\textwidth]{Figures/bg_fig1_ridership.pdf}
    \caption[EDSA Carousel ridership]{EDSA Carousel ridership. (a) Annual ridership
    grew 5.79\% from 63.02M in 2024 to 66.67M in 2025. (b) Daily ridership in 2025
    averaged 182{,}655 passengers, with a single-day peak of 321{,}186 in April 2025
    --- approximately 1.76$\times$ the daily mean \cite{DOTr2025Ridership}.}
    \label{fig:bg-ridership}
\end{figure}
%%

Figure~\ref{fig:bg-corridor-map} shows the EDSA Carousel Southbound route from Monumento to PITX, the corridor this study is grounded in, together with the other public transport modes (jeepney, MRT, LRT, tricycle, UV/FX) that intersect it at each major stop.

\begin{figure}[htbp]
    \centering
    \includegraphics[width=0.75\textwidth]{Figures/bg_fig3_edsa_corridor_map.pdf}
    \caption[EDSA Carousel corridor map]{EDSA Carousel Southbound route
    (Monumento to PITX) with intersecting public transport modes at each
    stop. Authors' illustration, adapted from the group's defense
    presentation.}
    \label{fig:bg-corridor-map}
\end{figure}

Despite the gradual increase in demand, the EDSA Bus Carousel still faces a lot of significant operational issues that affect the quality of bus routing. Firstly, the LTFRB relied on a tedious manual monitoring process by deploying personnel at each stop, creating decisions through empirical observation. \cite{EDSApolicy2023}. In relation to that, it is also clear that bus bunching has not been resolved; most of the time, their designated dispatch schedule is not followed \cite{Ollero2024EDSA}. There is an acknowledged lack of a structured monitoring framework for operations (e.g., dispatching, dwell time, bus bunching), in which adjustments are often done only after recurring issues surface \cite{Tiglao2025}. Moreover, mechanical failures or bus breakdowns are also prominent in Metro Manila, which contributes to a stressful environment, highlighting that poor vehicle condition and the constant threat of mid-route breakdowns are significant chronic stressors for Filipino bus drivers \cite{Chua2026}. Furthermore, environmental stochasticity severely destabilizes operations: empirical studies on Philippine expressways show that increasing rainfall intensity significantly reduces average traffic speed and free-flow capacity \cite{TSSP_Rain2018}. This rainfall-impact evidence is drawn from a 2018 study of the North Luzon Expressway rather than the EDSA Busway, and is used here only as contextual motivation that weather materially affects Philippine road-traffic operations; the weather-disturbance generator in this study (Section~3.2.6) does not adopt this study's specific speed-reduction percentages, and EDSA-specific travel-time behavior is independently calibrated through the GEH/RMSE procedure described in Section~3.2.3.

\begin{figure}[htbp]
    \centering
    \includegraphics[width=0.7\textwidth]{Figures/bg_fig2_rainfall_traffic.pdf}
    \caption[Rainfall impact on traffic flow]{Reduction in average speed and
    free-flow capacity relative to the clear-weather baseline across PAGASA
    rainfall categories. Capacity loss grows sharply at heavy intensities
    ($-17.4\%$), whereas mean speed degrades more moderately ($-7.4\%$), adapted from \cite{TSSP_Rain2018}.}
    \label{fig:bg-rainfall}
\end{figure}

As Figure~\ref{fig:bg-rainfall} shows, the reduction in
average speeds are about 5.34\% under light rain conditions, 6.3\% under moderate rain
conditions, and 7.4\% under heavy rain conditions. The average volume decreased by 2.92\%,
10.62\%, and 12.39\% while the capacities reduced by 3.67\%, 7.6\%, and 17.44\% under light,
medium, and heavy rain conditions, respectively \cite{TSSP_Rain2018}. Severe weather events often initiate localized flooding at key segments like Muñoz and Monumento Stations or cause severe degradation, forcing operational halts and emergency speed restrictions that further affect the busway’s performance  \cite{PhilstarTyphoon2024, PIA_Emergency2023, BusRepair2023}. Given these operational hurdles, the bus dispatching process of the EDSA Carousel can be modeled as a Vehicle Scheduling Problem (VSP).

Traditionally, the VSP has been addressed through static scheduling, which utilizes fixed timetables and headways derived from historical averages \cite{Ceder2007}. However, these rigid, rule-based methods often fail in stochastic environments as they cannot adapt to real-time disruptions, necessitating a shift toward more dynamic, autonomous scheduling frameworks that can respond to the unpredictability of the corridor \cite{Ollero2024EDSA}. This motivated a transition toward data-driven methodologies. Initial studies utilized traditional machine learning (ML) for demand forecasting and traffic prediction \cite{Wang2017}, but these models are primarily limited to predictive tasks and lack the active control mechanisms required for real-time decision-making in a closed-loop system. The subsequent wave introduced Single-Agent Reinforcement Learning (SARL), which enabled autonomous control, yet these frameworks frequently suffer from the `curse of dimensionality' as the joint state and action spaces grow exponentially with the number of vehicles, leading to centralized bottlenecks that hinder scalability in large transit networks \cite{Chen2024Survey}. To overcome these limitations, Multi-Agent Reinforcement Learning (MARL) has emerged as the state-of-the-art approach, offering a decentralized execution framework that enables scalable, collaborative control among multiple buses to mitigate complex issues like bus bunching \cite{Wang2020Holding}.

Although recent studies have demonstrated the potential of MARL-based dynamic scheduling to improve transit adaptability and overall reliability, a critical research gap remains. Existing MARL bus scheduling models are primarily evaluated in controlled simulation settings and are typically tested against isolated disruptions. No prior work has evaluated these models under the combined, simultaneous severe disruptions such as heavy rain, mechanical breakdowns, traffic variations, and sudden passenger surges that portray real-world corridors like EDSA. Therefore, this study aims to evaluate a MARL-based dynamic bus scheduling model under these non-ideal, stochastic traffic conditions to mitigate bus bunching, minimize passenger waiting times, and improve overall dispatch efficiency on the EDSA Carousel.

%%MARQUEZ END
\clearpage
\section{Review of Related Work}
This section surveys prior work on data-driven bus and traffic scheduling using a funnel structure. It starts from the broad transition from static scheduling toward adaptive and learning-based control, narrows to single-agent reinforcement learning and its known scalability limitations in large transit systems, narrows further to Multi-Agent Reinforcement Learning (MARL) and the assumptions commonly used in existing MARL bus-scheduling studies, and finally identifies the robustness gap addressed in this study.

\subsection{From Static Rules to Machine Learning}
\textbf{Static scheduling} refers to the offline determination of bus departure times, headways, fleet allocations, and stop assignments based on historical demand and travel-time averages, fixed in advance of operations and not adjusted in real time in response to observed disturbances \cite{Ceder2007}. Once deployed, such schedules cannot react to live disturbances such as demand surges, congestion, weather-induced delays, or vehicle breakdowns. The poor performance of fixed timetables under such conditions is well documented: even-headway and schedule-based control degrade rapidly once travel-time and demand variability exceed the narrow band assumed at design time, allowing small headway perturbations to amplify into bunching \cite{Daganzo2009}. Static schedules therefore remain mathematically inadequate for stochastic traffic environments, where the governing quantities are random variables rather than deterministic constants. Under the specific non-ideal conditions this study targets, the failure modes differ by control strategy. A fixed timetable has no feedback mechanism at all, so once bunching begins nothing in the schedule corrects it. A local reactive rule that holds a bus based only on the gap to the bus ahead can partially correct bunching under ordinary congestion, but has no way to respond to a breakdown, since it observes only the forward gap and not the enlarged gap a failed bus leaves behind it. A more globally aware reactive rule that accounts for both the forward and backward gap improves on this, but still follows a fixed, pre-specified rule rather than a learned response, so it cannot adapt its behavior to the heavier-tailed delays that severe weather introduces. These limitations motivated the transition toward more adaptive and data-driven scheduling methodologies.

Machine Learning (ML) techniques subsequently emerged as important tools for improving operational responsiveness in transportation systems through demand estimation, traffic forecasting, predictive dispatching, and adaptive holding strategies \cite{huang2019,patil2025}. Unlike static scheduling approaches, ML-assisted frameworks incorporate continuously updated operational information to improve short-term decision support under uncertain traffic conditions \cite{huang2019}. Recent applications to transit operations have demonstrated measurable operational gains. Barrera Hernandez et al. applied passenger-demand forecasting to support a heuristic dispatcher for bus operations in Bogotá, reporting reductions in passenger waiting time and operational cost relative to fixed timetables \cite{Barrera2025Optimization}. Similarly, Usman et al. utilized ML-based density estimation to improve real-time responsiveness in public transit systems \cite{Usman2025ML}. These studies demonstrate the growing role of ML-assisted methodologies in adaptive transit scheduling and operational planning.

Despite these advances, many ML-assisted transit-control frameworks still rely heavily on manually engineered intervention logic and externally designed dispatch heuristics \cite{alexandre2023}. In several predictive-control and heuristic-dispatch systems, forecasting components improve operational visibility, but the resulting operational responses are still governed by manually specified holding rules, threshold conditions, or predefined intervention strategies \cite{pang2019,verbich2021}. In the framework proposed by Barrera Hernandez et al., for example, the forecasting model informs the dispatcher but does not autonomously determine operational actions such as holding duration, dispatch timing, or stop-skipping behavior \cite{Barrera2025Optimization}. Consequently, while ML-assisted scheduling frameworks improve anticipatory planning and adaptive responsiveness, many existing implementations still stop short of establishing a fully integrated closed-loop control framework capable of autonomously optimizing sequential dispatch decisions over time.

\begin{table}[htbp]
\centering
\caption{Progression from static rules to reinforcement learning along the dimensions that determine real-time control capability.}
\label{tab:control_paradigms}

\footnotesize
\renewcommand{\arraystretch}{1.3}

\begin{tabularx}{\textwidth}{
>{\RaggedRight\arraybackslash}p{2.8cm}
>{\RaggedRight\arraybackslash}p{3.5cm}
>{\RaggedRight\arraybackslash}p{3.5cm}
>{\RaggedRight\arraybackslash}X
}
\toprule
\textbf{Paradigm} & 
\textbf{Adapts to real-time conditions?} & 
\textbf{Produces a dispatching action?} & 
\textbf{How the action is determined} \\
\midrule

Static / heuristic scheduling & 
No, fixed offline from historical averages & 
Yes, but actions are predetermined before operation & 
Pre-set timetables or manually designed heuristic rules that do not adapt to live corridor conditions \cite{Ceder2007,Daganzo2009} \\

Machine-learning assisted scheduling & 
Yes, forecasts update continuously using observed operational data & 
Partially; predictions support decisions, but actions are still governed externally & 
Forecasting models estimate demand or traffic conditions while separate heuristic dispatch rules determine the operational response \cite{Barrera2025Optimization,Usman2025ML,pang2019} \\

Reinforcement learning & 
Yes, policies continuously condition on live observed system states & 
Yes, the dispatching action is produced directly by the learned policy & 
The control policy learns through interaction with the environment to maximize long-term operational performance \cite{Momenikorbekandi2023Intelligent,Rodriguez2023Cooperative,Wangsun} \\

\bottomrule
\end{tabularx}
\end{table}

The decisive distinction is the final column: only reinforcement learning makes the dispatching action itself a learned output. This study instantiates the reinforcement-learning row as a parameter-sharing Double Deep Q-Network under centralized training with decentralized execution (Chapter~3).
\FloatBarrier
\subsection{Single-Agent Reinforcement Learning and Its Limitations}
Reinforcement Learning (RL) closes the prediction--control gap by learning a \textit{policy}: a mapping from observed system states to actions, optimized through trial-and-error interaction with the environment to maximize a long-horizon reward signal \cite{Momenikorbekandi2023Intelligent}. Unlike a forecaster, an RL agent does not just predict what will happen --- it decides what to do. The reward signal closes the loop: the agent's action changes the environment, the environment returns a new state and a scalar reward, and the policy is updated so that better actions become more likely the next time a similar state is encountered.

When RL is applied to a multi-vehicle problem such as bus-fleet control, the simplest formulation treats the whole fleet as one entity controlled by a single agent. This is \textbf{Single-Agent Reinforcement Learning (SARL)}: one centralized policy observes the full system state and emits a joint action covering every bus at once. SARL has been applied directly to bus holding control, including the STDH-DQN approach of Zhao et al.~\cite{Zhao2022STDH}, which uses a self-attention state encoder over spatial-temporal AVL features, and the more recent SA-DRL approach of Zhang and Zheng~\cite{Zhang2025SADRL}, which encodes categorical identifiers (vehicle ID, station ID, trip ID) as state features and reports favorable results against MARL baselines on a bidirectional timetabled line. SARL is therefore a real and competitive option in this space; the question is not whether it works but where it begins to strain.

SARL faces three scalability limitations in multi-bus control, each documented in the bus-scheduling and traffic-control literature. These limitations motivate the move to a multi-agent formulation in the next subsection; each has a corresponding design response that MARL provides.

\begin{enumerate}
\item \textbf{State-space dimensionality growth.} A SARL formulation of $N$ buses on $M$ stops requires a global state $s \in \mathbb{R}^{N \cdot d}$, where $d$ is the per-bus feature count (position, headway, load). For the EDSA Carousel sub-corridor studied here, with $M = 24$ stops and a fleet of $N \approx 12$--$30$ buses, $\dim(s)$ reaches roughly $96$--$240$. The function approximator must also learn approximate permutation invariance over buses, the policy should not depend on which arbitrary index a bus is assigned, a property not enforced by default in neural networks \cite{Wang2023MultiObj}.

\item \textbf{Joint action-space combinatorial growth.} A SARL controller emits a joint action $a = (a_1, \ldots, a_N) \in A^N$. With the hybrid action space adopted in this study ($|A_i| = 5 \times 2 = 10$, combining a 5-bin holding-strength discretization with a binary skip decision), a SARL controller managing $N = 20$ buses operates over a joint action space of $10^{20}$ combinations. Although deep RL policies do not enumerate these combinations directly --- they sample from a parameterized distribution over actions --- the credit-assignment burden remains: a single scalar reward must be attributed across all $N$ simultaneous component actions, and this attribution problem worsens as $N$ grows. MARL with shared policies factorizes the action selection into $N$ independent decisions of size $10$, each with its own attributable reward signal.

\item \textbf{Sample inefficiency under non-stationarity.} SARL treats the fleet as a single dynamical system; every demand change, breakdown, or weather event triggers re-adaptation of the entire policy. MARL with parameter sharing benefits from \textit{experience aggregation}: every action by every bus generates one training sample for the same shared network, so each bus's experience accelerates learning for all of them \cite{Christianos2021PS,Yu2022MAPPO}.
\end{enumerate}

Recent SARL-for-bus-control work that competes successfully against MARL baselines does so partly by hand-engineering categorical agent-identity features into the state representation \cite{Zhang2025SADRL}. This is informative because it suggests that agent identity remains an important representational issue in SARL formulations even at the high end of reported performance --- the dimensionality concern in item~(1) above is partially addressed by smuggling identity through input features rather than by a change in policy class. The fix is workable on a fixed corridor but does not transfer cleanly when fleet size or route topology changes.

We note a counter-position. Zhang and Zheng~\cite{Zhang2025SADRL} report SA-DRL matching or exceeding MARL baselines on bidirectional timetabled lines, attributing this to data-imbalance issues in MARL when some agents see few episodes. The EDSA Carousel sub-corridor studied here is operated as a high-frequency BRT service along a defined directional segment, closer to the single-direction loop-line assumptions where MARL has been shown to excel, which motivates the MARL choice here while acknowledging this caveat.

To situate the MARL literature reviewed in the next subsection within the broader ML and SARL landscape, Table~\ref{tab:ml_sarl_coverage} extends the paradigm comparison of Table~\ref{tab:control_paradigms} with a disturbance-coverage column, using the same D/S/T/W/B notation as Table~\ref{tab:marl_performance}.

\begin{table}[htbp]
\centering
\caption{Disturbance coverage across ML and SARL vehicle-scheduling studies, preceding the MARL-specific comparison in Table~\ref{tab:marl_performance}.}
\label{tab:ml_sarl_coverage}

\footnotesize
\renewcommand{\arraystretch}{1.3}

\begin{tabularx}{\textwidth}{
>{\RaggedRight\arraybackslash}p{2.6cm}
>{\RaggedRight\arraybackslash}p{2.0cm}
>{\RaggedRight\arraybackslash}X
>{\Centering\arraybackslash}p{1.6cm}
}

\toprule
\textbf{Paper} &
\textbf{Paradigm} &
\textbf{Method} &
\textbf{Disturbances} \\
\midrule

Wang et al.~\cite{Wang2017} &
ML (data-driven) &
Bus scheduling model incorporating time-dependent traffic and demand &
D \\

Barrera Hernandez et al.~\cite{Barrera2025Optimization} &
ML-assisted (heuristic dispatcher) &
Passenger-demand forecasting supporting a heuristic dispatcher &
D \\

Zhao et al.~\cite{Zhao2022STDH} &
SARL &
STDH-DQN; self-attention state encoder over spatial-temporal AVL features &
D, T \\

Zhang and Zheng~\cite{Zhang2025SADRL} &
SARL &
SA-DRL; categorical identity features (vehicle, station, trip ID) &
D, T \\

Verbich and El-Geneidy~\cite{verbich2021} &
Heuristic (non-MARL) &
Dynamic transit control under severe weather and vehicle breakdowns &
W, B \\

\bottomrule
\addlinespace[2pt]
\multicolumn{4}{@{}p{\textwidth}@{}}{\footnotesize D~=~stochastic demand, S~=~demand surges, T~=~traffic-speed perturbation, W~=~weather-induced heavy-tail delay, B~=~discrete vehicle breakdown. Disturbance coverage for the ML/SARL rows follows the characterization given by the RTC panel during the proposal review; unlike Table~\ref{tab:marl_performance}, these classifications have not been independently re-verified against each paper's full text.} \\
\end{tabularx}
\end{table}

The funnel is now complete: no ML or SARL study covers W or B, and among MARL studies (Table~\ref{tab:marl_performance}), only Verbich and El-Geneidy's heuristic controller~\cite{verbich2021} addresses both --- and it is explicitly non-MARL. Patil et al.~\cite{Patil2025Conformal} similarly validate weather-induced travel-time distributions but do not address bus control at all; their contribution to this study is the lognormal parameterization used by the weather-disturbance generator (Section~3.2.6), not a bus-control baseline. No prior study, ML, SARL, or MARL, combines W and B coverage with an actual MARL bus-scheduling controller, which is the specific gap this study fills.

\subsection{Multi-Agent Reinforcement Learning}

\begin{figure}[htbp]
    \centering
    \includegraphics[clip, trim=0.1cm 0.7cm 0.1cm 0.1cm,width=\textwidth]{Figures/rrl_fig1_sarl_vs_marl (3).pdf}
    \caption[SARL vs MARL architecture]{Comparison of single-agent and
    multi-agent formulations. (a) SARL: a centralized
    policy $\pi_\theta(s)$ consumes the concatenated global state
    $s \in \mathbb{R}^{N \cdot d}$ and outputs a joint action covering every bus.
    (b) MARL with parameter sharing: a single set of weights $\theta$ is shared
    across $N$ agents, each acting on its own local observation $o_i$. Authors' illustration.}
    \label{fig:sarl-vs-marl}
\end{figure}

In both panels of Figure~\ref{fig:sarl-vs-marl}, the per-bus state (denoted $s_{i,t}$ in the formal MDP notation of Section~\ref{subsec:state-space}, and shown in the figure as the local observation $o_i$) encodes the bus's current position, forward and backward headways, onboard load, and queue length at its current stop, as defined in full in Section~\ref{subsec:state-space}. The action $a_{i,t}$ is the holding-strength and stop-skipping decision the controller emits for that bus, defined in Section~\ref{subsec:action-space}. In the SARL panel (a), a single centralized network ingests all $N$ per-bus state vectors concatenated into one global state $s \in \mathbb{R}^{N \cdot d}$ and outputs all $N$ actions simultaneously; in the MARL panel (b), the same shared network weights $\theta$ instead process each bus's local state independently, so each agent acts on only its own observation rather than the concatenated global one.

Multi-Agent Reinforcement Learning (MARL) addresses the three limitations above by decomposing decision-making across multiple agents that share the environment. Each limitation has a direct design counterpart in MARL: state-space growth is contained by giving each agent only its own \textit{local} observation rather than the global state; action-space growth is factorized by having each agent select only its own action; and the sample-inefficiency problem is mitigated by sharing one network's parameters across all agents, so that every agent's experience updates the same policy. The problem is typically formalized as a Markov Game built on an underlying Markov Decision Process, with each agent maintaining its own local observation and acting under a shared or coordinated policy \cite{Li2023Departure,Sun2024Graph,Zuo2026AGV,Huang2025Joint,Bokade2023MARL}.

This decomposition introduces a new difficulty: \textit{partial observability}. No individual bus can see the full corridor state; it knows only its own headway, load, and immediate surroundings. A fully decentralized scheme in which each agent both trains and acts on its local observation alone would be unstable, because the other agents' simultaneously evolving policies make the environment non-stationary from any one agent's viewpoint --- the same local state can transition to different outcomes depending on what the other agents happen to be learning at that moment.

To address this, most modern formulations use \textbf{Centralized Training with Decentralized Execution (CTDE)}, illustrated in Figure~\ref{fig:ctde}. During training, agents are allowed to access system-wide information and optimize a shared objective; this stabilizes learning despite the non-stationarity introduced by simultaneously adapting agents. At deployment, each bus acts only on its own local observation, which is what makes the framework practically deployable on a real fleet without requiring a centralized real-time controller \cite{Cao2022Train,Wang2024MultiAGV,Nie2025CMRM,Xu2026Hierarchical}.



\begin{figure}[htbp]
    \centering
    \includegraphics[width=\textwidth]{Figures/rrl_fig2_ctde (2).pdf}
    \caption[Centralized training with decentralized execution]{(a) During training, all agents share one reward
    signal and one set of policy weights $\theta$, which are updated centrally
    using joint-state information. (b) At deployment, the weights are frozen and
    each bus runs an identical copy of $\pi_\theta$ on its own local observation
    $o_i$ alone, so no centralized real-time controller is required. Authors' illustration.}
    \label{fig:ctde}
\end{figure}

In the bus-scheduling setting, a parameter-sharing CTDE architecture is fleet-size invariant: the same trained policy applies regardless of how many buses are currently in service, because every bus is an instance of the same agent class acting on its own local headway, occupancy, and queue observations. Adding or removing a bus does not invalidate the trained model. Christianos et al.~\cite{Christianos2021PS} demonstrate that parameter sharing scales effectively to large agent populations and provides empirical support for the fleet-size invariance claim.

Recent progress has combined MARL with deep learning algorithms --- Deep Q-Networks (DQN), Proximal Policy Optimization (PPO) and its multi-agent variant MAPPO, and Graph Attention Networks (GAT) --- to handle high-dimensional observation spaces and complex interactions. These have been applied across vehicle scheduling, signal control, and logistics, generally reporting improvements in routing efficiency, congestion, and throughput in simulation \cite{Li2023Departure,Cao2022Train,Wang2024MultiAGV,Sun2024Graph,Zuo2026AGV,Huang2025Joint,Bokade2023MARL}.

\subsection{MARL Applied to Bus Scheduling}
\label{subsec:marl-applied}
Within bus scheduling specifically, the MARL literature has accumulated along a recognizable trajectory: first establishing that holding control could be learned at all, then enriching the action and objective space, and most recently beginning to examine policy behavior under disturbance. This trajectory is what the present study extends.

The foundational result is due to Wang and Sun~\cite{Wang2020Holding}, who showed that a cooperative MARL framework could learn an effective bus-holding policy on a single-line corridor and outperform classical headway-equalization rules under idealized stochastic demand. Chen et al.~\cite{Chen2016MARLBus} provided an earlier proof-of-concept in the same vein on a corridor with multi-agent reinforcement learning over local headway observations. These early results established the viability of the approach but were limited to the holding action alone, restricting the policy's leverage once bunching had already developed.

Subsequent work enriched the formulation along two axes. Rodriguez et al.~\cite{Rodriguez2023Cooperative} extended the action space by combining holding with stop-skipping in a cooperative DDQN framework on a simulated transit corridor with Poisson-distributed passenger arrivals and Gaussian inter-stop travel times, reporting reductions in average passenger waiting time relative to an even-headway holding baseline. Wang and Sun~\cite{Wang2023MultiObj} extended the objective space, scaling MARL to multi-line services with a shared corridor and jointly optimizing headway regularity against operational cost under controlled simulated demand and travel-time profiles.

Recent work has increasingly examined policy behavior under disturbance. Wang and Sun's IQNC-M~\cite{Wangsun} introduced implicit quantile networks and meta-learning for risk-sensitive bus holding, evaluating performance under demand surges and traffic-speed perturbations but not under weather-induced heavy-tailed delays or discrete vehicle breakdowns. In adjacent settings, Katzilieris et al.~\cite{Katzilieris2026MARL} integrated MADDPG-based traffic signal control with dynamic bus lane access management on a simulated arterial network with deterministic signal cycles; Liu et al.~\cite{Liu2023DRLHolding} reported DRL-based holding control under stochastic travel time and demand; and Shi et al.~\cite{Shi2022DistDRL} reported a distributed DRL bus control system in a connected environment with breakdown events modeled but without weather effects.

Across this trajectory, the simulation environments share three idealizing assumptions: passenger arrivals are stationary Poisson processes, inter-stop travel times are symmetric and well-behaved (Gaussian or similar), and the fleet experiences no discrete failure events. Reported performance gains are therefore conditional on those assumptions holding at deployment. The EDSA Carousel sub-corridor studied here violates each of them. Travel times are highly non-stationary, with tail behavior driven by tropical-cyclone and heavy-rainfall events that suspend operations entirely on the most severe days \cite{DOTr2020Suspension}; passenger demand swings sharply across peak intervals, with single-day ridership peaks roughly $1.8\times$ the daily average \cite{DOTr2025Ridership}; and tight headways mean that any individual vehicle breakdown propagates through the corridor as a downstream demand-absorption problem rather than as an isolated event. A controller whose reported gains rest on the idealizing assumptions above has no validated behavior on a corridor that violates all three at once.

\begin{table}[htbp]
\centering
\caption{Disturbance coverage across MARL bus-scheduling literature.}
\label{tab:marl_performance}

\footnotesize
\renewcommand{\arraystretch}{1.3}

\begin{tabularx}{\textwidth}{
>{\RaggedRight\arraybackslash}p{2.5cm}
>{\RaggedRight\arraybackslash}p{3.0cm}
>{\RaggedRight\arraybackslash}p{2.8cm}
>{\Centering\arraybackslash}p{2.2cm}
>{\RaggedRight\arraybackslash}X
}

\toprule
\textbf{Paper} &
\textbf{Algorithm} &
\textbf{Environment} &
\textbf{Disturbances} &
\textbf{Reported gain} \\
\midrule

Wang \& Sun (2020)~\cite{Wang2020Holding} &
Cooperative MARL holding &
Single-line, idealized &
D &
Bunching reduction vs.\ headway rule \\

Chen et al.\ (2016)~\cite{Chen2016MARLBus} &
Early MARL holding &
Corridor, idealized &
D &
Proof-of-concept gains \\

Rodriguez et al.\ (2023)~\cite{Rodriguez2023Cooperative} &
Cooperative DDQN (hybrid action) &
Chicago, mesoscopic &
T &
Wait time vs.\ headway-control baseline \\

Wang \& Sun (2023)~\cite{Wang2023MultiObj} &
Multi-objective MARL &
Shared corridor, idealized &
D &
Multi-line bunching reduction \\

Wang \& Sun (2023)~\cite{Wangsun} &
IQNC-M (distributional) &
Real-world bus services &
S, T &
Stable control under extreme events \\

Katzilieris et al.\ (2026)~\cite{Katzilieris2026MARL} &
MADDPG &
Athens arterial, SUMO &
T &
Throughput improvement \\

Liu et al.\ (2023)~\cite{Liu2023DRLHolding} &
DRL holding &
Synthetic single line &
D, T &
Wait time vs.\ no-control \\

Shi et al.\ (2022)~\cite{Shi2022DistDRL} &
Distributed DRL &
Connected environment &
D, B &
Integrated control gains \\

\textbf{This study} &
\textbf{DDQN + PS (CTDE)} &
\textbf{Calibrated EDSA SUMO} &
\textbf{S, T, W, B} &
\textbf{To be evaluated} \\

\bottomrule
\addlinespace[2pt]
\multicolumn{5}{@{}p{\textwidth}@{}}{\footnotesize D~=~stochastic demand, S~=~demand surges, T~=~traffic-speed perturbation, W~=~weather-induced heavy-tail delay, B~=~discrete vehicle breakdown.} \\
\end{tabularx}
\end{table}

\clearpage
Only Shi et al.~\cite{Shi2022DistDRL} carries a B (breakdown) entry in Table~\ref{tab:marl_performance}. Cao et al.~\cite{Cao2022Train}, which also models discrete vehicle failures, is deliberately excluded from this count: their MARL application is to \textit{train} rescheduling, not bus scheduling, so it does not belong in a table scoped to MARL bus-control literature. Verbich and El-Geneidy~\cite{verbich2021} likewise model breakdowns but use heuristic, non-MARL control (Table~\ref{tab:ml_sarl_coverage}), so they are excluded for the same reason. Among MARL bus-scheduling studies specifically, Shi et al. remains the only one to model discrete breakdowns.

Table~\ref{tab:marl_performance} summarizes what each study evaluated, what disturbances it modeled, and what it reported. No prior MARL bus paper has evaluated controller performance under a framework that incorporates all four disturbance classes documented in the EDSA corridor's operational profile such as demand surges, traffic-speed variability, heavy-tailed weather-induced travel-time delays, and discrete breakdowns in a calibrated real-corridor microsimulation. Of these four, weather-induced heavy-tailed delays and discrete breakdowns are the two that no prior MARL bus-scheduling study has modeled at all; their combination with demand surges and a calibrated traffic-speed baseline is the specific disturbance profile introduced here. The algorithmic choice this study adopts to address that combined-disturbance setting, DDQN with parameter sharing under CTDE is justified in detail in Chapter~3 (Section~\ref{subsec:proposed-algorithm}), alongside the rationale for not using actor-critic or distributional alternatives. The action space the agents operate on is correspondingly small and discrete: at each control event, an agent selects a holding strength from a five-bin discretization $\alpha \in \{0.0, 0.1, 0.2, 0.3, 0.4\}$ (multiplied by a maximum holding duration to yield seconds held) combined with a binary decision to skip the next stop, giving $|A_i| = 10$ actions per event. The full action-space rationale, including why discretization is preferred over a continuous holding parameter, is given in Chapter~3 (Section~\ref{subsec:action-space}).



Under the non-ideal evaluation regime the weather generator (W) \textbf{replaces} the traffic-speed generator (T) as the source of inter-stop travel-time variability rather than acting alongside it. Both are anchored to the same per-segment empirical mean $\mu$, and the weather coefficient of variation ($\eta \ge 0.3$) already spans and exceeds the $\pm 20\%$ range produced by traffic-speed scaling (Section~3.2.6). T therefore governs travel-time stochasticity under the ideal regime and W governs it under the non-ideal regime; the two are not composed on the same segment.



\subsection{The Disturbance Gap}
\label{subsec:disturbance-gap}
The idealized assumptions in Table~\ref{tab:marl_performance} limit how confidently the reported MARL gains can be expected to transfer to real urban corridors. Real corridors violate all three assumptions: demand exhibits non-stationary surges, weather-induced travel times are skewed and heavy-tailed, and individual buses break down stochastically. Wang and Sun themselves identify this as central, noting that existing MARL bus-control studies overlook performance under perturbations and that this becomes critical when transferring models toward real-world deployment \cite{Wangsun}. Their IQNC-M framework is among the first to address demand surges and traffic-speed perturbations in bus holding control, but weather-induced heavy-tailed delays and discrete vehicle breakdowns remain outside its tested scenario set.

Two considerations explain why the gap has persisted:

\begin{enumerate}
\item \textbf{Training instability under heavy-tailed returns.} Prior MARL bus-scheduling studies have typically used Poisson arrivals and Gaussian travel-time noise because these distributions generally produce more statistically stable training behavior than heavy-tailed disturbances. In reinforcement learning literature, heavy-tailed returns and stochastic gradients are known to introduce convergence difficulties and unstable optimization behavior during policy learning \cite{Henderson2018Matters,Agarwal2021Precipice}.
\clearpage
\item \textbf{Lack of temporally aligned anomaly data.} Historical severe-weather records are rarely aligned with operational bus data at the per-segment granularity needed for direct parameter estimation, making it difficult to fit an empirically grounded heavy-tailed disturbance distribution from corridor-specific data alone, especially given the short operational observation windows typically available. \textit{This study addresses the obstacle by adopting the lognormal travel-time parameterization that Patil et al.\ \cite{Patil2025Conformal} validated against INRIX freeway data via the Kolmogorov-Smirnov test, and then sweeping the disturbance intensity across a range of validated and extrapolated CV values rather than fitting a single point estimate. This characterizes controller behavior as a function of disturbance magnitude and avoids reliance on a corridor-specific severe-weather sample whose variance cannot be reliably resolved from a short observation window.}
\end{enumerate}

\subsection{Significance of Closing the Gap}
This study contributes the first characterization of MARL bus-scheduling performance across a swept range of disturbance intensities on a calibrated real-corridor simulation, with the disturbance set covering heavy-tailed weather-induced travel-time delays and discrete bus breakdowns that no prior MARL bus-scheduling study has modeled. Understanding policy performance under disturbance matters for two reasons. First, deploying MARL bus-scheduling models on the basis of ideal-simulator benchmarks alone leaves their real-world failure modes unknown until deployment, which is an unacceptable risk profile for a system that carries hundreds of thousands of daily passengers. Second, the relevant performance question for transit operators is not only average-case waiting time under quiet conditions but how a control policy degrades as a disruption intensifies, a characterization that prior MARL bus-scheduling work does not report. Establishing a reproducible stress-testing protocol pairing a calibrated corridor simulation with configurable demand, weather, and breakdown generators therefore contributes both an empirical characterization of performance under disturbance and a reusable evaluation framework for subsequent MARL work in this domain \cite{Li2023Departure,Wang2024MultiAGV,Che2024Recharging,Huang2025Joint,Bokade2023MARL}.
\chapter{Problem Statement}

\section{Rationale}
The public transportation system has long been a commodity that comes in the form of bus transport, train networks, and the like, with predetermined routes and fares \cite{Tironi2018}. In Metro Manila, traffic congestion has long been a problem stemming from insufficient road infrastructure and poor traffic mitigation policies, resulting in an inefficient public transport system \cite{Spatio2026}. One of the government's actions for an alternative transportation system that runs within the metropolitan area is the implementation of the EDSA Carousel in 2020. It is the only busway system in Metro Manila implemented by the government to have a designated lane, separated from other vehicles by concrete barriers \cite{Chua2026}. Over the years, the number of commuters using the EDSA Carousel has steadily risen, averaging approximately $1.83\times10^5$ daily passengers, with single-day peaks exceeding $3.2\times10^5$ during high-demand periods, reaching a new high of nearly 67 million in 2025, up from 63.02 million in 2024 
\cite{DOTr2025Ridership}. 


However, despite its growing demand, the system still faces several problems that can affect the quality of bus routing. To start, the LTFRB relied on a tedious manual monitoring process by deploying personnel at each stop \cite{EDSApolicy2023}. There are also certain portions where it operates with mixed traffic, causing delays among buses. It is also clear that bus bunching has not been resolved; most of the time, their designated dispatch schedule is not followed \cite{Ollero2024EDSA}. Moreover, the operation of the bus system is also largely influenced by external factors, such as weather and road maintenance. Sections of the EDSA Bus Carousel routes typically close during tropical cyclones to prevent potential damage from continuous rain and strong winds \cite{DOTr2020Suspension}. Given these operational issues, the bus dispatching process of the EDSA Bus Carousel can be modeled as a Vehicle Scheduling Problem (VSP), in which buses must be dispatched and coordinated efficiently despite these issues. By improving the scheduling, operational issues, such as bus bunching, irregular dispatch intervals, and passenger waiting times, will be reduced, thereby improving the overall bus dispatch efficiency.

Busway systems operate in a highly dynamic, unpredictable environment. In the past, Vehicle Scheduling Problems (VSP) relied on traditional scheduling techniques that focused on finding the optimal path by determining the shortest distance. Static scheduling is the pre-computed timetabling that fixes departure times before operation begins, with no mechanism to revise them in response to real-time conditions \cite{Ju2023Joint,Cao2022Train,Nie2025CMRM}. However, because transit environments are inherently stochastic, these static schedules perform poorly during daily operations, which causes severe operational issues such as bus bunching, node failures, and an increase in passenger wait times \cite{Daganzo2009,FEMbusbunching}.

To address the rigidity of static schedules, recent scheduling studies have shifted toward dynamic, AI-driven approaches, with Multi-Agent Reinforcement Learning (MARL) as the leading solution. Past studies have successfully utilized MARL to overcome the scalability limits of earlier systems by distributing the scheduling decisions across autonomous agents, improving the overall system's adaptability to real-time traffic \cite{Rodriguez2023Cooperative,Wang2023MultiObj,Wangsun}. However, a significant gap remains in how these state-of-the-art models are evaluated. Existing literature primarily validates the MARL scheduling models in a controlled, simulated environment. Previous evaluations rarely account for the heavy-tailed, asymmetric disruptions characteristic of actual transit operations, leaving their performance under non-ideal, real-world traffic anomalies largely untested. To the best of our knowledge, no existing MARL bus-scheduling model has been sufficiently tested against severe, simultaneous disruptions such as sudden passenger surges and weather-induced delays. To bridge this gap, this study develops and evaluates a MARL-based bus scheduling model within a stochastic traffic simulation subjected to non-ideal conditions.

This study develops and evaluates a MARL-based bus scheduling model on a calibrated stochastic simulation of the EDSA Carousel corridor, with non-ideal conditions operationalized as the simultaneous and independent injection of passenger demand surges, severe-weather travel-time skew, and Poisson-distributed bus breakdowns. The objective is to quantify how MARL performance degrades from its ideal-condition baseline as disturbance intensity rises, producing the first integrated robustness characterization for MARL bus scheduling under this specific combination of disturbances.

\section{Research Gap}
Existing Multi-Agent Reinforcement Learning frameworks for bus scheduling demonstrate measurable improvements in passenger waiting time and headway regularity, but these gains are validated almost exclusively in simulation environments that assume stationary passenger arrivals, symmetric Gaussian travel times, and a failure-free fleet. The performance of MARL bus-scheduling models has not been characterized when these idealizing assumptions are simultaneously violated by (i) stochastic passenger demand surges, (ii) heavy-tailed weather-induced travel-time delays, and (iii) discrete bus breakdowns. Without this characterization, it cannot be determined whether reported MARL gains persist, degrade gracefully, or collapse under realistic operating disturbances, which in turn blocks the transition of MARL bus scheduling from simulation to real urban transit deployment.

The weather-disturbance class (W) in particular was identified through the literature survey conducted earlier in this study (Section~\ref{subsec:marl-applied}), which found that no prior MARL bus-scheduling paper models heavy-tailed weather-induced travel-time delays (Table~\ref{tab:marl_performance}, column W). Its operational relevance to the EDSA corridor is established by the rainfall-driven reductions in average speed and free-flow capacity documented in Section~1.1 \cite{TSSP_Rain2018} and by the typhoon-related service suspensions recorded for the corridor \cite{DOTr2020Suspension}. The lognormal parameterization adopted for this disturbance class follows the Kolmogorov--Smirnov-validated form of Patil et al.~\cite{Patil2025Conformal}, introduced in this study to address the resulting lack of temporally aligned, corridor-specific anomaly data (Section~\ref{subsec:disturbance-gap}).



\section{Research Objectives}
\subsection{Main Objective}
To investigate the performance of a developed Multi-Agent Reinforcement Learning (MARL) scheduling model for dynamic bus scheduling under non-ideal traffic conditions on a calibrated simulation of the EDSA Carousel corridor.

\subsection{Specific Objectives}
\begin{enumerate}
   \item \textbf{Environment Construction:} To construct a stochastic traffic simulation environment representing a defined traffic area.

   \textbf{Expected Output 1.1:} A simulation model of a defined operational scope, validated against empirical traffic data to establish an environment that accurately approximates real-world traffic.

   \textbf{Expected Output 1.2:} An environment simulator capable of introducing environmental variability through adjustable parameters to simulate non-ideal operating conditions.

   \item \textbf{Algorithm Implementation:} To develop and implement a Multi-Agent Reinforcement Learning (MARL) approach for dynamic bus scheduling.

   \textbf{Expected Output 2.1:} A formulated agent architecture defining individual bus units as agents with discrete state and action space with a continuous reward function that punishes irregularities and passenger service delays.

   \item \textbf{Performance Evaluation Under Ideal and Non-Ideal Conditions:} To evaluate the scheduling performance of the MARL approach under both ideal operating conditions and non-ideal conditions.

   \textbf{Expected Output 3.1:} A quantitative and statistical assessment of MARL scheduling performance under ideal conditions, measuring passenger waiting time, travel time, and headway coefficient of variation, tested against baseline methods.

   \textbf{Expected Output 3.2:} A quantitative and statistical assessment of MARL scheduling performance under non-ideal conditions, measuring passenger waiting time, travel time, and headway coefficient of variation, reported as degradation relative to ideal-conditions performance.
\end{enumerate}


\section{Significance of the Study}
This study contributes both practical and scientific significance.
 
\textbf{Practical significance.} The EDSA Carousel carries on the order of $1.8\times10^5$ passengers daily, with single-day peaks exceeding $3.2\times10^5$ \cite{DOTr2025Ridership}. At this scale, even a small per-passenger improvement aggregates into a large operational effect: an illustrative 30-second reduction in mean passenger waiting time, applied across average daily ridership, corresponds to roughly $1.5\times10^3$ passenger-hours saved per day. (This figure is illustrative of the scale of impact rather than a performance target; actual savings depend on how uniformly the reduction is realized across riders and operating periods.) Beyond average-case waiting time, this study reports controller performance as a function of disturbance intensity, characterizing how a learned policy degrades as operating conditions worsen --- a perspective directly relevant to transit operators making service-reliability decisions during disruptions.

\textbf{Scientific significance.} The primary scientific contribution is empirical: this study produces the first characterization of whether reported MARL bus-scheduling gains persist, degrade gracefully, or collapse when the idealizing assumptions of prior work are violated simultaneously by demand surges, traffic-speed perturbations, weather-induced heavy-tailed delays, and discrete breakdowns. To produce this characterization reproducibly, the study contributes a stress-testing protocol for bus-scheduling reinforcement learning: a calibrated corridor model paired with four configurable stochastic anomaly generators. Because the protocol fixes the corridor, the disturbance generators, and the evaluation metrics, subsequent MARL work can be benchmarked on the same basis, allowing performance claims to be compared rather than reported in isolation. The study further reports the performance-degradation curve of MARL as a function of disturbance intensity, characterizing robustness rather than only average-case performance.

\section{Scope and Limitations}
 
\textbf{Scope.} This study develops and evaluates a MARL-based bus scheduling framework for the EDSA Carousel corridor. The framework is built on a calibrated SUMO microsimulation and runs over a single-day operational horizon; the simulation horizon, fleet size, and stop set are specified in Chapter 3. Each bus is modeled as an independent agent sharing one learned policy, observing only local conditions and acting at designated control stops. The MARL policy is benchmarked against three baselines — No Control, Forward Headway, and Even Headway — using passenger waiting time, total travel time, and headway coefficient of variation as metrics, over at least 30 Monte Carlo runs per regime. The same trained policy is then evaluated under two regimes, ideal and non-ideal, where non-ideal conditions are defined below.

For this study, \textbf{non-ideal conditions} are operating regimes in which one or more of four disturbance classes are active: (i) passenger demand surges, where boarding demand is scaled up above its baseline; (ii) traffic-speed perturbations, where corridor cruising speed varies around its baseline to represent everyday congestion friction; (iii) severe-weather travel-time delays, where inter-stop travel times are drawn from a heavy-tailed distribution to capture rain-induced slowdowns; and (iv) discrete bus breakdowns, where individual vehicles fail at random and leave service. The exact distributions and parameter ranges for each class are specified in Chapter~3. The severe-weather generator characterizes controller behavior across a range of travel-time variability magnitudes rather than across named meteorological categories; mapping specific intensity values to operational weather labels (e.g.\ ``moderate rain'' versus ``typhoon-grade'') is deferred to follow-on work that pairs the framework with a corridor-specific severe-weather travel-time dataset.
\clearpage
\textbf{Delimitations.} (a) Due to computational constraints, the simulation is restricted to a defined operational sub-segment of the EDSA Carousel corridor rather than the entire metropolitan road network, and minor feeder roads leading into the corridor are not modeled. Both restrictions are justified by the same structural fact: the EDSA Carousel operates on a physically separated, barrier-protected busway \cite{Chua2026}, so the agents' state and reward depend only on bus dynamics within the dedicated lane, specifically headways, dwell times, and onboard loads, none of which are directly observed by or computed from feeder-road traffic. Feeder roads affect the corridor only indirectly, through the passenger arrival rates they produce at each stop, and that effect is already captured by the calibrated per-stop demand distributions (Section~3.2.5) without needing to simulate the feeder network itself. Modeling feeder roads in SUMO would add computational cost without adding any new information the agents' observation or reward could use, since the sub-corridor restriction also preserves 1:1 empirical traffic volumes for GEH calibration without resorting to flow scaling; corresponding GEH calibration statistics are reported in Chapter~4. (b) MARL control actions are applied only at designated control stops, selected following criteria established in the literature (see Chapter~3). (c) The study is entirely simulation-based; no on-road deployment is performed. (d) Weather and breakdown effects are modeled as statistical distributions of their impact on travel time and fleet availability, not as physical or mechanical simulations of rain dynamics or vehicle subsystems. (e) Only the four disturbance classes defined above are considered; other disturbances (road accidents, signal failures, driver behavior variability) are out of scope. (f) Sim-to-real transfer is out of scope; the experimental claims are bounded to the calibrated SUMO microsimulation.
\chapter{Methods and Research Design}

This chapter describes the development, training, and evaluation of a Multi-Agent Reinforcement Learning (MARL) model for dynamic bus scheduling under non-ideal conditions. The goal is a reproducible pipeline from a calibrated stochastic microsimulation of the EDSA Carousel corridor to measurable performance results under both ideal and non-ideal operating conditions.


\section{Research Design}

The study uses a \textbf{controlled comparative experimental design} with two between-condition factors and three response variables. The factors are control strategy (No Control, Forward Headway, Even Headway, MARL) and operating condition (ideal, non-ideal). The response variables are mean passenger waiting time, mean total travel time, and headway coefficient of variation. The non-ideal condition is not a single setting but a range of disturbance intensities, so controller performance is characterized as a function of disturbance magnitude rather than evaluated at a single point.

Each combination is replicated $N \geq 30$ times using independent Monte Carlo seeds, so the Central Limit Theorem applies to the run means \cite{Wasserman2004}. The unit of analysis is one full simulated operating day on the defined EDSA Carousel sub-corridor. Random seeds and disturbance realizations are matched across control strategies within each disturbance level, so any performance gap between two controllers on the same run reflects the controller rather than the specific disturbance realization \cite{Henderson2018Matters,Agarwal2021Precipice}. Sampling is exhaustive within the simulation, so inclusion is defined at the disturbance level: only the modeled disturbances enter the non-ideal condition, and all others stay at their ideal values.

Throughout, disturbances play two distinct roles that must not be conflated. During \textbf{training}, the full range of stochastic disturbances is always active, the policy is trained under domain randomization \cite{Tang2024Curriculum} so that it learns to act under demand surges, traffic friction, heavy-tailed weather delays, and breakdowns. During \textbf{evaluation}, the same generators are used to define two test conditions (ideal and non-ideal) that probe the \textit{already-trained} policy on disturbance realizations it did not see during training but which are drawn from the same generative process. The ideal/non-ideal split is therefore an \textbf{evaluation factor, not a training switch}: the agent is never trained on a disturbance-free corridor. 
\pagebreak
\section{Methodology}

\subsection{Simulation Environment}

The study uses a two-stage simulation architecture. \textbf{SUMO} (Simulation of Urban Mobility)~\cite{Lopez2018SUMO} serves as the \textbf{calibration layer}, providing a microscopic representation of the EDSA Carousel corridor from which empirical parameters are extracted. A \textbf{custom Python environment built on the PettingZoo Agent Environment Cycle (AEC) API}~\cite{Terry2021PettingZoo} serves as the \textbf{training and Monte Carlo evaluation layer}, where MARL agents interact with a lightweight bus dynamics model parameterized by the distributions extracted from SUMO. SUMO is not used for training because its microscopic simulation of surrounding non-bus traffic is unnecessary for the agents' reward function, which depends only on bus state, and because stepping that traffic at every control event across the full Monte Carlo budget would be computationally prohibitive. SUMO is therefore used once for calibration, and all training and evaluation runs are carried out in the lightweight Python environment.

PettingZoo's AEC API is chosen over three alternatives. First, the single-agent Gymnasium interface~\cite{Towers2024Gymnasium} cannot natively represent multiple agents acting on independent events. Fitting multi-agent control into a single-agent API requires either treating the entire fleet as one composite agent, which breaks parameter sharing, or coordinating agents manually outside the API, which removes the benefit of using a standard interface. Second, PettingZoo's parallel API assumes all agents step simultaneously~\cite{Terry2021PettingZoo}, which does not match asynchronous bus control where only the bus arriving at a control stop has a decision to make. Third, a fully custom implementation was considered but rejected because it would sacrifice compatibility with established algorithm libraries and complicate reproducibility. AEC handles the asynchronous case by cycling through one agent at a time rather than stepping all agents in parallel~\cite{Terry2021PettingZoo}. This matches the approach used by Wang and Sun~\cite{Wangsun} and Rodriguez et al.~\cite{Rodriguez2023Cooperative}, both of whom implemented the same asynchronous logic in custom event-driven simulators. Adopting AEC preserves their event-driven dynamics while providing a standardized interface compatible with parameter-sharing algorithm libraries.

\begin{figure}[htbp]
    \centering
    \includegraphics[width=\textwidth]{Figures/fig_3_1_pipeline (2).pdf}
    \caption[Two-phase simulation pipeline]{The two-phase simulation architecture. Authors' illustration.}
    \label{fig:pipeline}
\end{figure}

SUMO is used for the calibration layer because it integrates with Python via the Traffic Control Interface (TraCI) and supports the GEH and RMSE validation protocols described later in this section. Ollero et al.~\cite{Ollero2024EDSA} previously developed an AIMSUN microsimulation of the EDSA Busway corridor; this study implements an equivalent network in SUMO to take advantage of its open-source TraCI interface, while adopting the same calibration targets.


A pure-Python alternative without SUMO calibration was also considered but rejected: it would lack the empirically grounded per-segment travel-time distributions that the GEH-and-RMSE-calibrated SUMO model provides. The chosen approach uses each tool where it is strongest: SUMO for calibrated microscopic realism during the one-time calibration phase, and Python for fast event-driven Monte Carlo during training and evaluation. This keeps the experiment tractable on accessible hardware, including a single cloud-notebook GPU runtime. Rodriguez et al.~\cite{Rodriguez2023Cooperative} and Wang and Sun~\cite{Wangsun} adopt the same two-layer approach, using event-based custom simulators with calibration parameters supplied from empirical data.



As shown in Figure~\ref{fig:pipeline}, the pipeline proceeds in two phases. First, SUMO is calibrated against empirical bus data, and the resulting per-segment travel time, demand, and dwell time distributions are saved. Second, the Python environment loads these distributions and runs a lightweight event-based bus simulator for MARL training and evaluation. Because the Python environment draws from distributions extracted from the calibrated SUMO model, its statistical behavior matches SUMO at the level of per-segment travel time, demand, and dwell time. It does not reproduce SUMO's microscopic vehicle-to-vehicle interactions, but those interactions affect the reward function only through their aggregate effect on inter-stop bus travel time, which the calibrated distributions already capture. Rodriguez et al.~\cite{Rodriguez2023Cooperative} make the same argument: they trained in an event-based mesoscopic simulator parameterized by empirical corridor data.

\subsection{Control-Stop Selection}
\label{subsec:control-stop-selection}

The study operates on a sub-corridor of the EDSA Carousel route. Within the chosen sub-corridor, only a subset of stops is designated as \textit{control stops} where MARL agents are permitted to act; the remainder are non-control stops that contribute stochastic dwell and delay behavior without agent intervention. Control stops are selected according to four criteria, following Rodriguez et al.~\cite{Rodriguez2023Cooperative} and the minimal-intervention principle of Wang and Sun~\cite{Wangsun}:

\begin{enumerate}
    \item \textbf{Origin terminal.} The route origin is always a control stop, since it is the natural point at which initial headways are set before buses enter the corridor.
    \item \textbf{Onset of high-demand segments.} The first stop of each high-demand segment is designated a control stop, so that the agent can act before demand-driven irregularity propagates downstream.
    \item \textbf{Low through-volume preference.} Stops with substantial through-passenger volume (many onboard riders not alighting) are \textit{avoided} as control stops, so that holding actions do not penalize large numbers of passengers already in transit.
    \item \textbf{No adjacent control hubs.} Where two high-demand hubs are adjacent, only the upstream hub is designated, so that holding penalties do not compound across two consecutive control actions.
\end{enumerate}

These criteria are data-driven: they specify \textit{how} stops are selected from demand and through-volume statistics, without assuming which stops a particular dataset will yield. Once the operational data are processed (Section~\ref{subsec:data-pipeline}) and the per-stop demand and through-volume profiles are known, the criteria are applied to produce the control-stop list for the study sub-corridor.

\subsection{Environment Model Validation}

Before MARL training begins, the SUMO calibration layer is statistically verified to approximate real-world bus operation along the corridor. Two industry-standard calibration metrics are used: the GEH statistic for bus volume on the dedicated corridor, and the Root Mean Square Error (RMSE) for bus speed trajectories. The calibration is restricted to the bus corridor itself, since the agents' state and reward depend only on bus dynamics; surrounding mixed-traffic flows do not enter the Python environment. This GEH/RMSE procedure calibrates EDSA-specific parameters directly from EDSA operational data and does not depend on the North Luzon Expressway rainfall-impact figures cited as motivating evidence in Section~1.1 \cite{TSSP_Rain2018}; that citation establishes only that weather materially affects Philippine road-traffic operations in general, not any EDSA-specific speed or capacity value used in this calibration.

\begin{figure}[htbp]
    \centering
    \includegraphics[width=\textwidth]{Figures/meth_fig_eo1_1_calibration (2).pdf}
    \caption[Illustrative SUMO calibration output]{Illustrative format of the
    SUMO calibration output. (a) Per-segment GEH statistic against the FHWA
    acceptance threshold of $\mathrm{GEH} < 5$; the corridor is accepted when at
    least 85\% of segments fall below the line. (b) Simulated bus-speed
    trajectory against the empirical observation used in the RMSE minimization
    step. Values shown are illustrative only. Authors' illustration.}
    \label{fig:calibration-illustrative}
\end{figure}

\subsubsection{Bus Volume Validation via GEH Statistic}

The GEH statistic, introduced for traffic-modelling acceptance testing in the UK Design Manual for Roads and Bridges \cite{DMRBV12}, measures the discrepancy between simulated and observed hourly bus volumes on individual corridor segments, illustrated in panel (a) of Figure~\ref{fig:calibration-illustrative}:

\begin{equation} GEH = \sqrt{\frac{2(M - C)^2}{M + C}} \label{eq:geh} \end{equation}

where $M$ is the hourly bus volume from the simulation and $C$ is the observed hourly bus count from empirical operational data. The environment is accepted when $GEH < 5$ for more than 85\% of corridor segments, following the FHWA Traffic Analysis Toolbox calibration criterion~\cite{FHWA2019TAT} as applied in recent SUMO calibration work by Patil et al.~\cite{Patil2025Conformal}.

\subsubsection{Speed Trajectory Calibration via RMSE}

Volume calibration alone does not guarantee that buses move realistically. RMSE evaluates how closely simulated bus speed trajectories match empirical observations, illustrated in panel (b) of Figure~\ref{fig:calibration-illustrative}:

\begin{equation} RMSE = \sqrt{\frac{\sum_{i=1}^{n}(x_{obs} - x_{model})^2}{n}} \label{eq:rmse} \end{equation}

where $x_{obs}$ is the observed bus speed and $x_{model}$ is the simulated bus speed \cite{Ollero2024EDSA}. To minimize RMSE, microscopic SUMO parameters like maximum desired speed and speed acceptance for the bus vehicle class are iteratively tuned to fit empirical speed profiles \cite{Ollero2024EDSA}. Once both the GEH threshold and RMSE minimization are achieved, the environment is treated as a validated representation of the operational scope.

\subsubsection{Real-Time Data Incorporation}

In a real-world deployment, the agent's state vector at each control event would be drawn from existing transit-data infrastructure: forward and backward headways $(h^-, \hat{h}^+)$ from Automatic Vehicle Location (AVL) feeds; onboard and waiting passenger counts from Automatic Passenger Counters (APC) and Automatic Fare Collection (AFC) terminals; and environmental flags from public weather APIs and incident-management systems \cite{Wangsun}. Within this study, the Python environment generates synthetic equivalents of these signals at each control event by querying the analytical bus model and combining outputs from the stochastic generators described below.

\subsection{Operating Conditions}

The study \textbf{evaluates} controllers under two distinct operating conditions. These conditions govern \textbf{evaluation only}; as detailed in the Training Protocol, the MARL policy is always trained with the full disturbance set active (domain randomization), and the two conditions below differ only in which disturbances are switched on when measuring the frozen, already-trained policy. The symbols used to define these conditions ($\eta$, $\sigma_d$, $\sigma_s$, and the breakdown rate $\lambda$) are collected in Table~\ref{tab:notation}.

\textbf{Ideal conditions.} The corridor operates under the empirical baseline from the operational dataset, with no additional disturbance. The weather disturbance generator is disabled ($\eta = 0$), the breakdown generator is disabled ($\lambda = 0$), and the demand and traffic-speed scaling factors are at their baseline means ($\sigma_d = 0$, $\sigma_s = 0$). Each bus therefore encounters the empirical mean travel time and demand profile for its segment and time of day. The Python environment still injects normal day-to-day variability through the calibrated per-segment distributions, so the ideal condition is \textit{stochastic but not disturbed} which comparable to a normal weekday with no incidents, severe weather, or surge events. This serves as the control condition for Stage A \textit{evaluation}; it is not a training setting, since training always exposes the policy to the full disturbance range.

\textbf{Non-ideal conditions.} One or more of the four stochastic generators described in the next section is activated. The full non-ideal evaluation (Stage B) sweeps the weather disturbance intensity $\eta$ across five values, with the breakdown generator active alongside throughout. Demand scaling remains active under both conditions to preserve realistic passenger-arrival variability. The traffic-speed generator governs inter-stop travel-time variability under the ideal condition; under the non-ideal condition, the heavier-tailed weather generator replaces it as the source of travel-time stochasticity, since its coefficient of variation $\eta \geq 0.3$ already spans and exceeds the variability range that the traffic-speed scaling produces (Section~\ref{subsec:stochastic-vars}). The two conditions therefore differ in which generator governs travel-time stochasticity and in whether the discrete breakdown generator is enabled.

Both conditions draw inter-stop travel time from distributions anchored to the same per-segment empirical baseline mean $\mu$ extracted during SUMO calibration; only the shape of variability changes between conditions. The non-ideal condition recovers the ideal condition as $\eta \to 0$, because the lognormal collapses to a point mass at $\mu$ when its coefficient of variation is zero. Performance differences between the two conditions therefore reflect a change in the variability profile of travel time, not a change in its mean or in the underlying network geometry.

Throughout this chapter, \textit{ideal conditions} and \textit{non-ideal conditions} refer to these operating states of the simulated environment. \textit{Baseline controllers} refers separately to the three non-MARL control strategies (No Control, Forward Headway, Even Headway) against which the MARL policy is benchmarked. A condition is a state of the world; a controller is a choice of algorithm.

\begin{table}[htbp]
\centering
\caption{Simulation parameter summary: fixed, swept/variable, and derived parameters.}
\label{tab:sim-parameters}
\small
\renewcommand{\arraystretch}{1.25}

\begin{tabularx}{\textwidth}{
>{\RaggedRight\arraybackslash}p{0.24\textwidth}
>{\RaggedRight\arraybackslash}p{0.09\textwidth}
X
>{\RaggedRight\arraybackslash}p{0.20\textwidth}
}

\toprule
\textbf{Parameter} & \textbf{Symbol} & \textbf{Value / Range} & \textbf{Source} \\
\midrule
\multicolumn{4}{l}{\textit{Fixed simulation parameters}} \\
\midrule

Simulation horizon & --- & Single simulated operating day (\%TODO-VAL: exact start/end hours) & Section~3.1 \\
Total stop count (sub-corridor) & $M$ & 24 & Introduction, Section~1.2.2 \\
Fleet size (active buses) & $N$ & $\approx 12$--$30$ & Introduction, Section~1.2.2 \\
Control stop count & --- & \%TODO-VAL: determined by criteria in Section~3.2.2 once dataset is processed & Section~3.2.2 \\
Scheduled headway & $H_0$ & \%TODO-VAL: from DOTr schedule records & DOTr records \\
Bus passenger capacity & --- & \%TODO-VAL: from DOTr fleet specification & DOTr records \\
Maximum holding duration & $\Delta T$ & \%TODO-VAL: to be set during implementation & Section~3.2.7 (Action Space) \\
Holding action bins & $\Omega$ & $\{0.0, 0.1, 0.2, 0.3, 0.4\}$ & Section~3.2.7 (Action Space) \\
Action space size per agent & $|A_i|$ & 10 (5 hold $\times$ 2 skip) & Section~3.2.7 (Action Space) \\
Monte Carlo runs per cell & $N_{\text{runs}}$ & $\geq 30$ & Section~3.1 \\
Discount / event-based discount & $\gamma$, $\beta$ & \%TODO-VAL: tuned during implementation & Section~3.2.7 \\

\midrule
\multicolumn{4}{l}{\textit{Swept / variable parameters}} \\
\midrule

Weather disturbance intensity & $\eta$ & $\{0.0, 0.3, 0.6, 1.0, 1.3\}$ & Section~3.2.6 \\
Demand scaling std.\ dev.\ (clip) & $\sigma_d$ & Clip $[1, 3]$ & Section~3.2.6 \\
Traffic-speed scaling std.\ dev.\ (clip) & $\sigma_s$ & Clip $[0.8, 1.2]$ & Section~3.2.6 \\
Breakdown rate & $\lambda$ & \%TODO-VAL: calibrated during implementation (events/hour) & Section~3.2.6 \\

\midrule
\multicolumn{4}{l}{\textit{Derived parameters (from SUMO calibration)}} \\
\midrule

Baseline inter-stop travel time & $\mu$ & Per-segment, per-time-of-day bin, \%TODO-DATA & SafeTravelPH dataset \\
Baseline travel-time std.\ dev.\ & $\sigma$ & Per-segment, per-time-of-day bin, \%TODO-DATA & SafeTravelPH dataset \\
Baseline coefficient of variation & $CV_0$ & $\sigma/\mu$ per segment, \%TODO-DATA & Computed from dataset \\
Lognormal shape parameter & $\sigma_{ln}$ & $\sqrt{\ln(\eta^2+1)}$ (Eq.~\ref{eq:sigma_ln}) & Method of moments \\
Lognormal location parameter & $\mu_{ln}$ & $\ln(\mu) - \sigma_{ln}^2/2$ (Eq.~\ref{eq:mu_ln}) & Method of moments \\

\bottomrule
\end{tabularx}
\end{table}

Table~\ref{tab:sim-parameters} therefore serves as the single reference point for every parameter used across this chapter. Parameters marked \%TODO-VAL are to be confirmed during the implementation phase upon receipt of the operational dataset and DOTr schedule records; parameters marked \%TODO-DATA will be computed during the SUMO calibration phase described in Section~3.2.3. The stop count ($M=24$) and fleet-size range ($N \approx 12$--$30$) are carried over from the state-space dimensionality discussion in Section~1.2.2 and are not new values introduced here.

\subsection{Data Processing}

The MARL framework is trained and evaluated on a single stream of input data pre-processed before training: corridor operational data used to calibrate SUMO and to derive baseline parameters for the stochastic generators. A second stream, the synthetic state observations generated by the Python environment at runtime, is produced on-the-fly during training and consumed once per control event.

\subsubsection{Required Datasets}

\begin{itemize}

\item \textbf{Corridor bus operational data.} A per-trip record of EDSA Carousel bus operation along the study sub-corridor, collected over a continuous observation window of at least two weeks. The required fields are GPS-tracked vehicle location, boarding and alighting events, passenger occupancy, operating speed, and \textit{dwell time} at each stop, the dwell time being the interval a bus spends stationary at a stop serving passengers, measured from the moment the doors open to the moment they close and the bus is ready to depart, exclusive of any holding time subsequently imposed by the controller. These records yield the empirical distributions of bus cruising speed, inter-stop travel time, and demand under ideal operating conditions, used both to calibrate SUMO and to define the baseline operating point of the stochastic generators. The baseline operating point for this study is established from a crowdsourced operational record collected from the EDSA Busway during July 2023 through the SafeTravelPH mobile application.

%A crowdsourced operational dataset of this form, for example, data collected from the EDSA Busway via the SafeTravelPH mobile application, provides a representative model of the required structure. Should bus-volume coverage prove insufficient for the GEH validation step, supplementary records may be requested from MMDA, DOTr, or the involved bus operators under the Freedom of Information process.

\end{itemize}

Table~\ref{tab:field-mapping} maps each required raw field to the parameter derived from it and the MARL component that parameter feeds into, connecting the data requirements above to the disturbance generators (Section~\ref{subsec:stochastic-vars}), the control-stop selection criteria (Section~\ref{subsec:control-stop-selection}), and the agent observation vector (Section~\ref{subsec:state-space}). The mapping reflects the study's design intent, not properties of a processed dataset; specific statistics remain \%TODO-DATA pending dataset acquisition.

\begin{table}[htbp]
\centering
\caption{Mapping of required raw dataset fields to derived parameters and their role in the MARL formulation.}
\label{tab:field-mapping}
\small
\renewcommand{\arraystretch}{1.25}

\begin{tabularx}{\textwidth}{
>{\RaggedRight\arraybackslash}p{0.19\textwidth}
>{\RaggedRight\arraybackslash}p{0.28\textwidth}
>{\RaggedRight\arraybackslash}X
}

\toprule
\textbf{Raw Field} & \textbf{Derived Parameter} & \textbf{MARL Component} \\
\midrule

GPS-tracked vehicle location &
Per-segment inter-stop travel time distribution ($\mu$, $\sigma$) &
SUMO speed calibration (RMSE, Section~\ref{subsec:stochastic-vars}); anchors the traffic-speed generator ($\sigma_s$) and the weather generator's baseline mean ($\mu$) \\

Boarding events &
Per-(stop, time-of-day) baseline demand rate &
Demand surge generator baseline ($\sigma_d$); waiting-passenger-count feature in the agent's local observation \\

Alighting events &
Through-passenger volume per stop &
Control-stop selection criterion 3 (low through-volume preference, Section~\ref{subsec:control-stop-selection}) \\

Passenger occupancy &
Per-segment load profile &
Onboard-passenger-count feature in the agent's local observation; dwell-time estimation \\

Operating speed &
Per-segment baseline cruising speed &
SUMO bus-volume calibration (GEH); traffic-speed generator baseline ($\sigma_s$) \\

Dwell time &
Per-(stop, time-of-day) dwell time distribution &
Event-driven bus model, used between control events to advance the simulation clock \\

\bottomrule
\end{tabularx}
\end{table}

Severe-weather conditions are not estimated from operational data in this study but are injected as a controlled experimental variable, with disturbance magnitudes anchored to validated literature values rather than to a corridor-specific severe-weather sample.

A short empirical observation window cannot reliably resolve the variance of a heavy-tailed severe-weather distribution at per-segment granularity, and a swept-disturbance design characterizes controller performance across a range of conditions rather than at a single uncertain calibration point. Because the disturbance conditions are random by construction, they are not drawn from any single fixed dataset: each Monte Carlo run generates a fresh, independently seeded realization of demand, traffic-speed, weather, and breakdown disturbances. The single empirical dataset is required only to fix the \textit{baseline} (ideal-condition) operating point against which these random disturbances are layered; the held-out validation split (Stage 3 below) confirms that the calibrated baseline generalizes beyond the records used to tune it. The full disturbance parameterization is described in Section~\ref{subsec:stochastic-vars}.

\subsubsection{Data Pre-Processing Pipeline}
\label{subsec:data-pipeline}

Pre-processing proceeds in three stages.

\textit{Stage 1: Cleaning.} Trip records with missing GPS coordinates, missing timestamps, negative inter-stop times, or trips that fail integrity checks (for example, a later stop served before an earlier one) are dropped. Remaining records are normalized to a common time zone. Because the operational observation window is short (on the order of two weeks), the data cannot support separate weekend or holiday demand profiles --- a window of this length yields too few weekend days and typically no holidays to estimate distinct calendars reliably. Records are therefore filtered to regular weekdays and binned by time-of-day, which is the calendar dimension that drives demand and travel-time variation within the modeled operating day. The restriction to weekday operation is recorded as a scope limitation.

\textit{Stage 2: Empirical distribution extraction.} For each (segment, time-of-day) bin, the cleaned operational data yield the mean $\mu$ and standard deviation $\sigma$ of inter-stop travel time under baseline operating conditions, producing the baseline coefficient of variation $CV_0 = \sigma / \mu$ used to anchor the disturbance generators. Boarding and alighting counts per (stop, time-of-day) bin yield the baseline demand profile. These statistics are serialized to disk as parameter files that the Python environment loads at initialization.

\textit{Stage 3: Train/validation split for calibration.} The calibration data are split chronologically into a calibration set, used to tune SUMO parameters until the GEH and RMSE targets are met, and a held-out validation set, used to confirm that the calibrated environment generalizes beyond the tuning data.

The MARL agents are not trained on these raw datasets directly. They are trained on synthetic state vectors $s_{i,t}$ produced by the Python environment at each control event. Each $s_{i,t}$ is assembled at runtime by querying the analytical bus model and is consumed once by the policy network before being discarded. There is no persistent training dataset of state-action transitions on disk other than the standard replay buffer used by the learning algorithm.

\subsection{Stochastic Disturbance Generators}
\label{subsec:stochastic-vars}

\paragraph{Disturbance Classes and Independence}

This study distinguishes five disturbance classes, denoted D, S, T, W, and B:

\begin{itemize}
\item \textbf{Stochastic demand (D):} the baseline, always-present day-to-day randomness in passenger arrivals, drawn from the calibrated per-(stop, time-of-day) demand distributions (Section~\ref{subsec:data-pipeline}). D is not a disturbance layered on top of a deterministic baseline; it \textit{is} the baseline stochastic environment, present in every run regardless of which other generators are active.
\item \textbf{Demand surge (S):} an episode-level multiplicative scaling factor, with standard deviation $\sigma_d$ (Table~\ref{tab:notation}), that amplifies baseline boarding rates above their empirical mean. S is the controlled experimental variable; D is always present, and S is what is added on top of it. Setting $\sigma_d = 0$ removes the surge and leaves only baseline demand variability (D).
\item \textbf{Traffic-speed perturbation (T):} an episode-level scaling of corridor cruising speed, with standard deviation $\sigma_s$, representing everyday congestion friction. T governs inter-stop travel-time variability under ideal conditions.
\item \textbf{Weather-induced delay (W):} a per-segment travel-time distribution with coefficient of variation $\eta$, drawn from a right-skewed lognormal rather than the Gaussian-based scaling used by T. W replaces T as the source of travel-time stochasticity once $\eta > 0$ (Section~\ref{subsec:stochastic-vars}).
\item \textbf{Discrete bus breakdown (B):} a Poisson-distributed discrete event, with rate $\lambda$, that permanently removes one bus from the active agent set for the remainder of the simulated day.
\end{itemize}

The four generators that produce S, T, W, and B are injected independently: no causal chain links them within the simulation. A breakdown event (B) does not trigger a demand surge (S) or a weather delay (W), and a weather event does not induce a mechanical failure. In practice, some real-world disturbances co-occur causally --- for example, heavy rain may both slow buses (W) and concentrate passengers at covered stops (S) --- but this study treats each generator as an independent factor. This design choice isolates the individual and combined effect of each disturbance class on controller performance and allows the single-disturbance ablation (Section~\ref{subsec:evaluation}) to attribute degradation unambiguously to a specific class.

Four stochastic generators inject variability into the Python environment. Generators (i) and (ii) follow the perturbation framework of Wang and Sun~\cite{Wangsun}; the weather generator's heavy-tailed lognormal formulation follows Patil et al.~\cite{Patil2025Conformal}; the breakdown generator follows the rescheduling formulation of Cao et al.~\cite{Cao2022Train}. Table~\ref{tab:notation} collects the symbols used across this section and the MARL formulation that follows.


\begin{table}[htbp]
\centering
\caption{Notation used in the stochastic generators and the MARL formulation.}
\label{tab:notation}

\small
\renewcommand{\arraystretch}{1.25}

\begin{tabularx}{\textwidth}{
>{\RaggedRight\arraybackslash}p{0.14\textwidth}
>{\RaggedRight\arraybackslash}X
>{\RaggedRight\arraybackslash}p{0.18\textwidth}
}

\toprule
\textbf{Symbol} & \textbf{Definition} & \textbf{Domain / units} \\
\midrule

$\eta$ &
Weather disturbance intensity; directly sets the coefficient of variation of disturbed inter-stop travel time ($CV=\eta$). $\eta=0$ disables the weather generator. &
dimensionless \\

$\sigma_d$ &
Standard deviation of the per-episode passenger-demand scaling factor. $\sigma_d=0$ fixes demand at its baseline mean. &
dimensionless \\

$\sigma_s$ &
Standard deviation of the per-episode traffic-speed scaling factor. $\sigma_s=0$ fixes cruising speed at its baseline mean. &
dimensionless \\

$\lambda$ &
Rate parameter of the Poisson process governing bus breakdowns. $\lambda=0$ disables breakdowns. &
events/hour \\

$\mu$ &
Baseline (empirical) mean inter-stop travel time for a given segment and time-of-day bin. &
seconds \\

$\sigma$ &
Baseline standard deviation of inter-stop travel time. &
seconds \\

$CV_0$ &
Baseline coefficient of variation, $CV_0=\sigma/\mu$, used to anchor the disturbance generators. &
dimensionless \\

$\mu_{ln},\sigma_{ln}$ &
Location and shape parameters of the lognormal weather-delay distribution (method-of-moments). &
--- \\

$h^-,\hat{h}^+$ &
Forward headway (gap to bus ahead) and estimated backward headway (gap to bus behind). &
seconds \\

$H_0$ &
Desired (scheduled) headway. &
seconds \\

$\alpha_{i,t}$ &
Holding-strength action of agent $i$, $\alpha_{i,t}\in\Omega=\{0.0,0.1,0.2,0.3,0.4\}$, scaled by $\Delta T$. &
dimensionless \\

$\Delta T$ &
Maximum holding duration that $\alpha_{i,t}$ scales. &
seconds \\

$\gamma,\beta$ &
Discount factor and event-based discount rate ($e^{-\beta t}$ replaces $\gamma$ for random inter-event times). &
dimensionless \\

$\theta$ &
Weights (parameters) of the online neural network (shared policy). &
dimensionless \\

$\theta^-$ &
Weights (parameters) of the target neural network in DDQN. &
dimensionless \\

$N$ &
Number of bus agents (also used for the Monte Carlo run count, $N\geq30$, where noted). &
count \\

\bottomrule
\end{tabularx}
\end{table}

\subsubsection{Random Seed Control}

All four generators draw from a single seeded pseudo-random number stream maintained by the Python NumPy library. At the start of each Monte Carlo run, a master seed is set in the simulation driver script, which initializes NumPy's BitGenerator. For each (disturbance level, control strategy) cell, the list of $N \geq 30$ seeds is fixed in advance and stored as a configuration file in the project repository. Within a disturbance level, the same seed list is reused across all four control strategies, so the demand trace, disturbance realization, and breakdown timing for run $k$ are identical regardless of which controller is active. This matching enables paired statistical comparisons: any difference in performance metrics between two controllers on run $k$ is attributable to the controller rather than to the disturbance realization \cite{Henderson2018Matters}. Seeds across runs are independent, and the choice $N \geq 30$ ensures that the Central Limit Theorem applies to the run-mean of each response variable. %For settings with heavy-tailed disturbances, bootstrap-based confidence intervals are reported rather than relying solely on the normal approximation at $N = 30$ \cite{Agarwal2021Precipice}.

\subsubsection{Passenger Demand}

The baseline empirical transit demand is perturbed each episode by a scaling factor sampled from $\mathcal{N}(1, \sigma_d^2)$ and clipped to $[1, 3]$, following the general Gaussian-clipped demand-scaling mechanism of Wang and Sun~\cite{Wangsun}, though this study adopts a narrower clip than their $[1, 10]$ range. The asymmetric clip focuses the test on demand surges rather than symmetric variation, since demand drops produce lightly loaded conditions that do not stress-test the controller. The upper bound of 3, corresponding to roughly a tripling of baseline boarding rates, is this study's own choice (\%TODO-VAL: revisit against Wang and Sun's wider range during implementation) rather than a value drawn from prior work. Sampling occurs at the start of each simulation run, producing varied demand profiles across episodes. In implementation, the scaling factor $f_d \sim \mathcal{N}(1, \sigma_d^2)$ is sampled once per episode at initialization and applied uniformly to every per-stop, per-time-of-day arrival rate for the duration of that simulated operating day, so all stops experience the same proportional demand shift within a single run while the shift itself varies across runs.

\subsubsection{Traffic Delays}

Mean cruising speed between stops is adjusted dynamically using a scaling factor drawn from $\mathcal{N}(1, \sigma_s^2)$, clipped to $[0.8, 1.2]$, representing typical daily congestion friction \cite{Wangsun}. In implementation, the speed scaling factor $f_s \sim \mathcal{N}(1, \sigma_s^2)$ is sampled once per episode and applied to the bus's mean cruising speed on every inter-stop segment traversal during that day, producing a uniformly slower or faster corridor for that run without segment-level variation beyond the calibrated baseline. This generator provides the baseline stochastic variability in inter-stop travel time when the weather generator is inactive. When the weather generator is active ($\eta > 0$), the heavier-tailed lognormal described below replaces this scaling as the source of travel-time variability: the lognormal coefficient of variation $\eta \geq 0.3$ used in the non-ideal sweep already exceeds the $\pm 20\%$ range that traffic-speed scaling produces, so the two generators do not compose on the same segment. The traffic-speed generator's empirical role persists through the per-segment baseline mean $\mu$ to which both generators are anchored.

\subsubsection{Weather-Induced Anomalies}

Severe-weather disturbances are injected as a controlled experimental variable rather than estimated from a fixed empirical sample.

The controlled variable in this generator is the coefficient of variation of inter-stop travel time, not a meteorological quantity. Severe weather is the motivating phenomenon, but the variable that enters the simulator and that the controller observes is the travel-time distribution itself. The lognormal is chosen to model the shape that any such heavy-tailed disruption produces on per-segment travel time, regardless of its meteorological label. In implementation, when $\eta > 0$ a fresh travel-time sample $T \sim \text{LogNormal}(\mu_{ln}, \sigma_{ln})$ is drawn independently for each bus at each inter-stop segment traversal during the episode, replacing the traffic-speed generator's output for that traversal; the lognormal parameters $\mu_{ln}$ and $\sigma_{ln}$ are computed from the segment's empirical mean $\mu$ and the swept $\eta$ via Equations~\eqref{eq:sigma_ln}--\eqref{eq:mu_ln}. The disturbance intensity sweep is therefore read as a span of travel-time variability magnitudes, not as a sweep across named weather categories.

Inter-stop travel time under disturbance is sampled from a Lognormal distribution, which captures the right-skewed, heavy-tailed delays observed in urban traffic under inclement weather \cite{Patil2025Conformal}. The lognormal is chosen over a Gaussian distribution, which is symmetric and cannot represent heavy right tails, and over Gamma or Weibull distributions, which can produce heavy tails but lack the direct empirical validation against urban traffic data that Patil et al.\ provide for the lognormal form. The distribution's mean is the baseline inter-stop travel time $\mu$ from the operational data, and its variability is set by a disturbance intensity parameter $\eta$, which directly specifies the coefficient of variation:

\begin{equation} CV = \eta \label{eq:cv_eta} \end{equation}

The lognormal shape and scale parameters follow the method-of-moments derivation:

\begin{equation} \sigma_{ln} = \sqrt{\ln(\eta^2 + 1)} \label{eq:sigma_ln} \end{equation}

\begin{equation} \mu_{ln} = \ln(\mu) - \frac{\sigma_{ln}^2}{2} \label{eq:mu_ln} \end{equation}

Travel time is drawn as $T \sim \text{LogNormal}(\mu_{ln}, \sigma_{ln})$. Patil et al.~\cite{Patil2025Conformal} tested this parameterization by generating SUMO-simulated travel times under the same CV-driven lognormal recipe --- with time windows and mean travel times anchored to INRIX historical data for an urban arterial corridor, not a freeway --- and confirming via the Kolmogorov-Smirnov test that the simulated distribution matches the assumed log-normal shape, reporting a close fit at the highest variability level they tested ($KS = 0.036$, $p = 0.94$ at $CV = 1.0$).

The disturbance intensity $\eta$ is swept across five values, $\eta \in \{0.0, 0.3, 0.6, 1.0, 1.3\}$, chosen to span and slightly exceed the empirically validated CV range of Patil et al.~\cite{Patil2025Conformal}, who evaluate scenarios at $CV \in \{0.1, 0.3, 0.5, 0.7, 1.0\}$. The basis for each swept value is as follows. $\eta = 0$ disables the generator and recovers the ideal condition. $\eta \in \{0.3, 0.6\}$ fall inside the validated range, representing moderate travel-time variability. $\eta = 1.0$ corresponds to the upper end of Patil et al.'s validated range and is adopted here as the nominal severe-disturbance operating point. $\eta = 1.3$ extends beyond the validated range as an extrapolated stress test; results at this level are interpreted as a probe of controller behavior under extreme conditions rather than as a calibrated scenario. Table~\ref{tab:eta-basis} summarizes this reasoning.

\begin{table}[htbp]
\centering
\caption{Basis for each swept weather-disturbance intensity value.}
\label{tab:eta-basis}
\small
\renewcommand{\arraystretch}{1.25}

\begin{tabularx}{\textwidth}{
>{\Centering\arraybackslash}p{0.08\textwidth}
X
}

\toprule
\textbf{$\eta$} & \textbf{Basis} \\
\midrule

0.0 & Generator off --- recovers the ideal baseline \\
0.3 & Inside Patil et al.'s validated range --- moderate variability \\
0.6 & Inside Patil et al.'s validated range --- moderate variability \\
1.0 & Top of Patil et al.'s validated range --- nominal severe operating point \\
1.3 & Beyond the validated range --- deliberate extrapolated stress test \\

\bottomrule
\end{tabularx}
\end{table}

The mapping from $\eta$ to a specific named weather severity (e.g.\ ``moderate rain'' versus ``typhoon-grade'') is not established by Patil et al., whose CV values are not labeled by weather category. The $\eta$ sweep should therefore be read as a span of travel-time variability magnitudes rather than as calibrated meteorological intensities. Reporting controller performance across this sweep characterizes degradation as a function of disturbance magnitude, which is methodologically stronger than evaluating at a single point.

\subsubsection{Bus Breakdowns}

Unlike the perturbations above, a bus breakdown is a \textit{discrete event} that removes a bus from active simulation \cite{Cao2022Train,Shi2022DistDRL}. Breakdowns are triggered at random times sampled from a Poisson process with a configurable rate $\lambda$ (Table~\ref{tab:notation}). In implementation, at each discrete simulation timestep of length $dt$, a Bernoulli trial with probability $\lambda \cdot dt$ is evaluated independently for each active bus; a success removes that bus from the active agent set for the remainder of the simulated day. When a breakdown occurs at bus $b_k$, $b_k$ is removed from the active agent set, leaving an enlarged headway between the trailing bus $b_{k-1}$ and the leading bus $b_{k+1}$. The passengers that $b_k$ would have served accumulate at the stops ahead of $b_{k-1}$, so the trailing bus inherits this residual demand as it advances.

A breakdown does not override any agent's policy; it changes the \textbf{state} the surviving agents observe and the \textbf{reward landscape} they face. The removal of $b_k$ is encoded in the downstream incident/breakdown flag of the trailing agent's observation (Section~\ref{subsec:state-space}), and the inherited residual demand sharply raises the headway-irregularity and dwell-time penalties that the trailing agent's reward already captures. Under these conditions, the stop-skipping action, which is part of the agent's normal action space (Section~\ref{subsec:action-space}) becomes strongly advantageous. Skipping a heavily loaded downstream stop lets $b_{k-1}$ shed onboard passengers (alighting permitted, boarding suppressed) and close the gap to $b_{k+1}$ before bunching sets in. The environment therefore \textit{induces} the recovery behavior through the state and reward signal, and the agent must \textit{learn} to select it; the action is never injected by the environment. This keeps the breakdown response consistent with decentralized execution and the Markov-game formulation, in which every action originates from $\pi_\theta$.

The Poisson assumption is a tractable approximation for the single-day simulation horizon, over which vehicle age, maintenance schedule, and mileage since last service can be treated as constant.


\subsection{MARL Architecture and MDP Formulation}

This subsection builds the formal learning framework in two steps. It first introduces the Bellman recursion that underlies value-based reinforcement learning, then specifies the Markov game tuple, the per-agent observation and action spaces, the reward-function design, and the proposed learning algorithm. The justification for choosing a multi-agent formulation over a single-agent alternative is presented in Chapter~1.

\subsubsection{Bellman Foundation}

At the core of reinforcement learning is the dynamic programming principle established by Bellman~\cite{Bellman1957}. The agent's objective is to maximize the expected sum of discounted future rewards. The state-action value function $Q^\pi(s,a)$, denoting the expected return from taking action $a$ in state $s$ and following policy $\pi$ thereafter, satisfies the Bellman optimality equation:

\begin{equation}
Q^*(s, a) = \mathbb{E} \left[ r(s, a) + \gamma \max_{a'} Q^*(s', a') \mid s, a \right]
\label{eq:bellman}
\end{equation}

where $s'$ is the state reached after applying action $a$, $r(s,a)$ is the immediate reward, and $\gamma \in [0, 1)$ is the discount factor. This equation is relevant here for two reasons. First, it shows that a long-horizon control policy can be learned by recursively backing up local one-step rewards, making it computationally feasible to evaluate the downstream effects of a holding or skipping decision without enumerating every future stop on the route. Second, in the asynchronous bus-control setting, the time between two consecutive control events is itself random. Following Bradtke and Duff~\cite{Bradtke1995}, the discount factor in Eq.~\eqref{eq:bellman} is replaced by an event-based discount $e^{-\beta t}$ that scales with the elapsed time between events. Each agent's learning target is therefore an event-adjusted version of Eq.~\eqref{eq:bellman}, not a fixed-tick MDP target. By convention throughout this chapter, $\gamma$ denotes the nominal discount factor as it appears in the canonical Bellman and DDQN equations, while $e^{-\beta t}$ is the event-based discount actually applied in the asynchronous setting; wherever $\gamma$ appears in a value-function equation, it is instantiated by $e^{-\beta t}$ at runtime.


\subsubsection{Markov Game Formulation}

Following Littman~\cite{Littman1994}, the bus fleet control problem is formulated as a multi-agent extension of a Markov Decision Process, a Markov game defined by the tuple $G = \langle N, S, A, P, R, \gamma \rangle$:

\begin{itemize}
\item $N$: number of bus agents.
\item $S$: joint state space.
\item $A$: joint action space.
\item $P$: state-transition function $P(s_{t+1}\mid s_t, a_t)$.
\item $R$: reward function.
\item $\gamma \in [0, 1)$: discount factor.
\end{itemize}

Because bus control is asynchronous and event-driven, agents act only upon arrival at a control stop. The formulation maintains a per-agent local transition $P_i : S_i \times A_i \rightarrow S_i$, with $S_i$ the local observation of agent $i$ \cite{Wangsun,Rodriguez2023Cooperative}.

\subsubsection{Agents ($N$)}

The decision-making agents are the individual bus units $i \in N$ operating along the corridor. To avoid combinatorial blow-up and to keep the system invariant to fleet size, all agents share a single neural-network policy under a \textbf{parameter-sharing} architecture and act on their own local observations at deployment \cite{Wangsun,Rodriguez2023Cooperative,Christianos2021PS}. Parameter sharing means that one set of network weights serves every bus; each bus computes its own action from its own observation, but learning updates from any bus improve the policy for all.

\subsubsection{State Space and Local Observations ($S_i$)}
\label{subsec:state-space}

The local observation $s_{i,t}$ at agent $i$ on arrival at a control stop is a vector composed of:

\begin{itemize}
\item \textbf{Spatial location:} index of the current control stop.
\item \textbf{System regularity:} forward headway $h^-$ and estimated backward headway $\hat{h}^+$, both available from AVL feeds in deployment \cite{Wangsun}.
\item \textbf{Passenger demand:} onboard passenger count and estimated number of passengers waiting at the stop, available from APC/AFC systems \cite{Wangsun}.
\item \textbf{Environmental flags:} encoded indicators for the current disturbance intensity and any active downstream incident or breakdown.
\end{itemize}

\begin{table}[htbp]
\centering
\caption{Agent observation vector: features, symbols, and data sources.}
\label{tab:observation-features}
\small
\renewcommand{\arraystretch}{1.25}

\begin{tabularx}{\textwidth}{
>{\RaggedRight\arraybackslash}p{0.24\textwidth}
>{\RaggedRight\arraybackslash}p{0.11\textwidth}
>{\RaggedRight\arraybackslash}X
>{\RaggedRight\arraybackslash}X
}

\toprule
\textbf{Feature} & \textbf{Symbol} & \textbf{Deployment Source} & \textbf{Simulation Source} \\
\midrule

Control stop index & --- & Route map (static) & Hardcoded stop list \\
Forward headway & $h^-$ & AVL feed & Event-driven bus model \\
Estimated backward headway & $\hat{h}^+$ & AVL feed & Event-driven bus model \\
Onboard passenger count & --- & APC system & Running tally in bus model \\
Waiting passenger count at stop & --- & AFC terminal / platform sensors & Stop queue in bus model \\
Disturbance intensity flag & $\eta$ (encoded) & Weather API / public incident alert & Active generator parameter \\
Downstream incident / breakdown flag & $b$ (binary) & Incident management system & Breakdown generator output \\

\bottomrule
\end{tabularx}
\end{table}

As Table~\ref{tab:observation-features} shows, in simulation all observation features are generated synthetically by the Python environment at each control event by querying the analytical bus model and the active stochastic generators; no real sensor data is consumed during training or evaluation.

\subsubsection{Action Space ($A_i$)}
\label{subsec:action-space}

The agent's action $a_{i,t} = \pi_\theta(s_{i,t})$ combines two discrete components, following Rodriguez et al.~\cite{Rodriguez2023Cooperative}:

\begin{itemize}
\item \textbf{Holding control:} a holding-strength parameter $\alpha_{i,t} \in \Omega = \{0.0, 0.1, 0.2, 0.3, 0.4\}$, multiplied by a maximum holding duration $\Delta T$ to yield seconds held; $\alpha_{i,t} = 0$ means no holding.
\item \textbf{Stop-skipping:} a binary decision to bypass the next stop, triggered when the bus is severely delayed and near capacity, allowing headway recovery. Skipping is disallowed at the origin terminal and is not permitted if the immediately preceding trip already skipped that stop, so no passenger is denied service twice in a row.
\end{itemize}

The full action set is the Cartesian product of these two components: $|A_i| = 5 \times 2 = 10$ discrete actions per control event, allowing the agent to select a holding strength and a skip decision independently at each control event. This is a broader action space than Rodriguez et al.~\cite{Rodriguez2023Cooperative}, whose combined holding-and-skipping controller (DDQN-HA) instead selects among six \textit{mutually exclusive} actions: five holding strengths $\Omega = \{0.0, 0.1, 0.2, 0.3, 0.4\}$ (with $\omega = 0$ already covering the no-holding case) plus a single skip action. The discretized holding-strength set $\Omega$ adopted here matches theirs exactly. A continuous holding parameter $\alpha \in [0, 1]$ was considered, following Wang and Sun~\cite{Wangsun}, but rejected for two reasons. First, continuous actions require actor-critic algorithms, whose training instability compounds across the swept-disturbance evaluation budget. Second, real driver compliance with holding instructions is itself imperfect: Rodriguez et al.~\cite{Rodriguez2023Cooperative} model non-compliant drivers as departing after only 60--80\% of the instructed holding time, so continuous precision in $\alpha$ is not meaningful at deployment.

\subsubsection{Reward Function ($R$) and Objective}
\label{subsec:reward}

The infinite-horizon objective of agent $i$ is the expected cumulative event-discounted return

\begin{equation}
R_i = \mathbb{E}\!\left[\sum_{k=0}^{\infty} e^{-\beta t_k}\, r_{i, t+k}\right],
\end{equation}

where $t_k$ is the elapsed time to the $k$-th future control event, consistent with the event-based discount introduced earlier.

The \textbf{structure} of the per-event reward $r_{i, t+k}$ --- its component terms and the operational priorities they encode --- is defined in this chapter. Its \textbf{exact weighting and sensitivity analysis} are a research objective of this study (Specific Objective 2, Expected Output 2.1), to be finalized during implementation rather than inherited from prior work. This separation is deliberate: the reward \textit{form} is grounded here, while the tuning of its coefficients is treated as an empirical question. The reward formulation must balance three competing operational priorities: (i) headway regularity, since bunching is the primary failure mode of bus corridors; (ii) passenger waiting time at served stops, which is the most direct passenger-experience metric; and (iii) a penalty term that discourages degenerate skipping behavior, since an unconstrained skip action could otherwise reduce headway variance by refusing to serve passengers. Existing MARL bus-control reward formulations span a range of trade-offs among these priorities, from headway-coefficient-of-variation forms for holding-only action spaces \cite{Wangsun,Wang2023MultiObj} to passenger-time forms for combined holding-and-skipping action spaces \cite{Rodriguez2023Cooperative}. This study establishes the overall reward structure for the hybrid action space by defining the three reward components and their additive formulation, while treating the corresponding weighting coefficients, together with their sensitivity analysis, as the implementation-phase deliverable under Expected Output 2.1. Although the component structure is fixed at this stage, the coefficients remain as placeholders to be determined during implementation through experimental evaluation.

The reward is computed independently for each agent at every control event rather than as a shared team-level objective. Accordingly, $r_{i,t}$ represents the reward assigned to agent $i$ and is stored as that agent's transition in the shared replay buffer (Training and Execution Protocol, step 4). Each bus is therefore evaluated based on the consequences of its own action, even though all agents learn from a common shared network. Since the reward is derived from locally observable quantities already contained in the agent's observation, particularly the forward and backward headway measurements, the resulting signal remains aligned with corridor-wide service regularity without requiring a centralized reward computation during execution.

The overall reward function is expressed as the weighted sum of three penalty terms:

\begin{equation}
r_{i,t+k} = -w_1 \cdot (\text{headway-irregularity term}) - w_2 \cdot (\text{waiting-time term}) - w_3 \cdot (\text{skip-degeneracy term}),
\label{eq:reward-form}
\end{equation}

where $w_1$, $w_2$, and $w_3$ denote the weighting coefficients (\%TODO-VAL) to be determined through the Expected Output 2.1 sensitivity analysis. Each component is formulated as a non-positive penalty, allowing the agent to maximize its cumulative return in Eq.~\eqref{eq:bellman} by minimizing headway irregularity, passenger waiting time, and unnecessary stop-skipping behavior. Consequently, this chapter establishes the reward formulation and its optimization objective, while the specific mathematical expressions and coefficient values are reserved for the implementation and evaluation phase.

\subsubsection{Proposed Learning Algorithm}
\label{subsec:proposed-algorithm}

The fully discrete action space ($|A_i| = 10$) makes value-based learning a natural fit. This study adopts \textbf{Double Deep Q-Network (DDQN)~\cite{vanHasselt2016DDQN} with parameter sharing under centralized training and decentralized execution (CTDE)}. CTDE means the policy is trained with access to information from all agents but each bus acts on its own local observation alone at deployment, so the trained system does not require a centralized real-time controller. The choice follows Rodriguez et al.~\cite{Rodriguez2023Cooperative}, who used the same DDQN-based framework for cooperative holding-and-skipping control on a comparable corridor.

DDQN is chosen over three categories of alternatives. \textit{Plain DQN}~\cite{Mnih2015DQN} overestimates Q-values in noisy environments \cite{vanHasselt2016DDQN,Henderson2018Matters}, which the high-variance disturbance regime in this study would amplify. DDQN's double estimator addresses this by decoupling action selection from evaluation \cite{vanHasselt2016DDQN}. \textit{Actor-critic methods such as MAPPO and MADDPG} are compatible with continuous action spaces but unnecessary given the discrete action-space decision above; they would also introduce training instability that compounds across the swept-disturbance evaluation budget. \textit{Distributional methods such as Wang and Sun's IQNC-M} \cite{Wangsun} model the full return distribution rather than its expectation and offer strong performance on this problem class, but their meta-learning training procedure adds implementation risk beyond the scope of a thesis-scale project; IQNC-M is noted as a comparison target for future work.

Two enhancements from the recent bus-control literature will be tested during implementation: extending agent awareness to include neighboring-trip information in the observation and reward, as in Rodriguez et al.'s DDQN-HA variant \cite{Rodriguez2023Cooperative}, and credit assignment via inductive graph learning, as in Wang and Sun's CAAC framework \cite{Wangsun}. The final architecture will be selected based on training stability and validation-set passenger waiting time after a fixed training budget.


\subsubsection{Double Deep Q-Network Foundation}
\label{subsec:ddqn-foundation}

The Bellman optimality equation defines the value $Q^*(s,a)$ that an ideal controller would assign to each state-action pair, but it cannot be solved exactly for a state space as large as the bus-control problem, where the observation is a continuous vector of headways, loads, and queues. Two approximations bridge this gap. $Q$-learning~\cite{Bellman1957,Littman1994} replaces the exact solution with an iterative update that nudges a stored estimate $Q(s,a)$ toward the one-step Bellman target $r + \gamma \max_{a'} Q(s',a')$ after each transition. The Deep $Q$-Network (DQN)~\cite{Mnih2015DQN} replaces the lookup table with a neural network $Q(s,a;\theta)$, so that the value function generalizes across the continuous state vectors the bus environment produces. DQN trains $\theta$ by minimizing the temporal-difference loss between the network's prediction and a target computed from a periodically frozen copy of the network, with transitions drawn from an experience replay buffer to decorrelate consecutive samples.

DQN has a known issue: it systematically overestimates action values~\cite{vanHasselt2016DDQN}. The cause is the single $\max$ operator in the target. Because the same network both selects the next action ($\arg\max_{a'}$) and evaluates it ($\max_{a'} Q$), noise in value estimates is amplified the $\max$ preferentially picks whichever action happens to be overvalued, and that inflated value propagates into the target. In a low-variance environment this bias is manageable, but the swept-disturbance regime in this study generates heavy-tailed returns that feed the overestimation mechanism directly.

The Double Deep $Q$-Network (DDQN)~\cite{vanHasselt2016DDQN} removes this bias by decoupling action selection from action evaluation. The online network (weights $\theta$) selects the next action, but a separate target network (weights $\theta^-$) evaluates that action's value:

\begin{equation}
y^{\text{DDQN}} = r(s,a) + \gamma\, Q\!\left(s',\, \arg\max_{a'} Q(s', a'; \theta);\; \theta^-\right),
\label{eq:ddqn_target}
\end{equation}

For the overestimation to persist, both networks would have to overvalue the same action simultaneously, which is less likely than a single network doing so alone. The network is trained by minimizing the squared temporal-difference error:

\begin{equation}
\mathcal{L}(\theta) = \mathbb{E}_{(s,a,r,s') \sim \mathcal{D}} \left[ \left( y^{\text{DDQN}} - Q(s,a;\theta) \right)^2 \right],
\label{eq:ddqn_loss}
\end{equation}

where $\mathcal{D}$ is the experience replay buffer and the target weights $\theta^-$ are held fixed for a number of updates (or slowly tracked toward $\theta$) so that the regression target does not shift at every step.

Three properties make DDQN appropriate for this study. First, the discrete action space ($|A_i| = 10$) is the setting value-based methods handle well: the $\arg\max$ in Eq.~\eqref{eq:ddqn_target} is a finite enumeration over ten options. Second, the double estimator addresses the overestimation that high-variance disturbance returns would otherwise amplify. Third, DDQN inherits DQN's experience-replay and target-network machinery, which provides sufficient training stability without the auxiliary mechanisms that actor-critic methods require; this point is discussed further in Section~\ref{sec:challenges}. The event-based discount $e^{-\beta t}$ replaces the fixed $\gamma$ in Eq.~\eqref{eq:ddqn_target} when the target is computed in the asynchronous setting, so the DDQN target inherits the same event-adjustment as the underlying Bellman recursion. The combination of DDQN with parameter sharing and the CTDE scheme is specified in the Proposed Learning Algorithm subsection.

\subsubsection{Training and Execution Protocol}

The learning process follows the standard reinforcement-learning feedback loop under the asynchronous, event-driven bus-control setting and the CTDE scheme described above, illustrated in Figure~\ref{fig:aec-training}. The loop, repeated at every control event during training, proceeds as follows:

\begin{enumerate}
\item \textbf{Observe.} When bus agent $i$ arrives at a control stop, the Python environment assembles its local observation $s_{i,t}$ (Section~\ref{subsec:state-space}) from the analytical bus model and the active stochastic generators.
\item \textbf{Act.} The shared policy network selects an action $a_{i,t}$ (holding strength and skip decision) from $s_{i,t}$ using an $\epsilon$-greedy rule over the estimated $Q$-values during training, and greedily ($\arg\max_a Q$) at evaluation.
\item \textbf{Transition and reward.} The environment applies $a_{i,t}$, advances the simulation to the next control event, and returns the event-discounted reward $r_{i,t}$ (Section~\ref{subsec:reward}).
\item \textbf{Store.} The transition $(s_{i,t}, a_{i,t}, r_{i,t}, s_{i,t'})$, where $t'$ is the agent's next control event, is written to a shared experience replay buffer common to all agents under parameter sharing.
\item \textbf{Update.} Mini-batches sampled from the replay buffer update the shared online network by minimizing the DDQN temporal-difference loss against a slowly-updated target network; experience from every bus updates the single shared parameter set.
\end{enumerate}

This loop is the feedback mechanism that distinguishes the learned controller from the fixed-rule baselines, which compute a holding time directly from the current headways with no update step.

The protocol is split into a training phase and an execution (evaluation) phase:

\begin{itemize}
\item \textbf{Centralized training phase.} The loop above is run for a fixed episode budget with the stochastic disturbance generators active, so that the policy encounters demand surges, traffic-speed friction, weather delays, and breakdowns during training. A policy that never sees a breakdown or heavy-tailed delay during learning has no basis for responding to one at deployment. This follows the domain-randomization approach of Tang et al.~\cite{Tang2024Curriculum}, who randomize demand and schedule variables during training to obtain policies that generalize to real-world traffic variability. Training uses system-wide information where it stabilizes learning (e.g.\ a shared reward and shared buffer), and experience from all buses contributes to one shared network.

\item \textbf{Decentralized execution phase.} The trained network's weights are frozen. During Monte Carlo evaluation (Stages A and B), each bus acts independently and greedily on its own local observation $s_{i,t}$ alone, with no centralized controller and no further weight updates. The fixed policy is evaluated across both the ideal condition (Stage A) and the swept non-ideal condition (Stage B), so that its behavior is measured on disturbance realizations it did not see during training but which are drawn from the same generative process.
\end{itemize}

\begin{figure}[htbp]
    \centering
    \includegraphics[width=\textwidth]{Figures/fig_3_2_aec_training.pdf}
    \caption[Agent--environment cycle (training)]{The asynchronous agent--environment
    cycle during training. Authors' illustration.}
    \label{fig:aec-training}
\end{figure}
 

Because heavy-tailed returns can destabilize value-based learning if injected without safeguards \cite{Henderson2018Matters,Agarwal2021Precipice}, training under disturbance is paired with the stabilizers discussed in Section~\ref{sec:challenges}: DDQN's double estimator, reward clipping, and slow target-network updates. If training under the full disturbance set proves unstable at the highest intensities despite these measures, a curriculum that ramps disturbance intensity over the course of training~\cite{Tang2024Curriculum} is retained as a fallback.

\clearpage

\subsection{Baseline Controllers for Comparison}

The MARL policy is benchmarked against three baseline controllers: \textbf{No Control (NC)}, \textbf{Forward Headway (FH)}, and \textbf{Even Headway (EH)}. These are the standard benchmark set in the MARL bus-control literature; both Wang and Sun \cite{Wangsun} and Rodriguez et al.~\cite{Rodriguez2023Cooperative} use this trio. Using the same baselines makes the results directly comparable to published work.

\subsubsection{No Control (NC)}

No holding or skipping is applied. Each bus dwells only for the time required to serve boarding and alighting passengers, then departs immediately. NC serves as the lower bound of the comparison: it isolates the natural corridor behavior without any regularization. If MARL cannot beat NC, the controller is doing more harm than good. NC also provides the reference point for measuring the severity of bus bunching. Under non-ideal conditions, NC has no corrective mechanism whatsoever, so demand surges, weather-induced delays, and breakdowns are expected to compound directly into bunching with no attenuation.

\subsubsection{Forward Headway (FH)}

The classical holding rule introduced by Daganzo~\cite{Daganzo2009}. When a bus arrives at a control stop, its holding time is:

\begin{equation}
d_{FH} = \max\!\left\{0,\, \bar{d} + g\,(H_0 - h^-)\right\}
\label{eq:fh}
\end{equation}

where $H_0$ is the desired (scheduled) headway, $h^-$ is the forward headway (time gap to the bus ahead), $\bar{d}$ is the average slack delay at equilibrium, and $g > 0$ is a control gain parameter. If a bus is running closer to the bus ahead than scheduled ($h^- < H_0$), it holds for an additional period proportional to the gap deficit; if it is already at or beyond the desired spacing ($h^- \geq H_0$), it does not hold. Parameters $\bar{d}$ and $g$ follow Daganzo's original values. FH represents local, single-agent reactive control: each bus reacts only to information about the bus immediately ahead. Because Eq.~\eqref{eq:fh} conditions only on $h^-$, FH is expected to partially correct demand- and traffic-induced bunching but to respond poorly to a breakdown: it has no observation of the enlarged gap a failed bus leaves for the bus behind it, so the trailing bus's holding decision is uninformed by the breakdown until the irregularity has already propagated back to it.

\subsubsection{Even Headway (EH)}

A more globally aware rule that holds buses to equalize the forward and backward gaps simultaneously, following Rodriguez et al.~\cite{Rodriguez2023Cooperative}:

\begin{equation}
T_{EH} = \max\!\left\{\min\!\left[\frac{\hat{h}^+ - h^-}{2},\; 0.4\,H_0\right],\, 0\right\}
\label{eq:eh}
\end{equation}

where $h^-$ is the forward headway, $\hat{h}^+$ is the estimated backward headway (time gap to the bus behind), and $\frac{\hat{h}^+ - h^-}{2}$ is the holding time that would split the difference evenly between the two gaps. The hold is capped at $0.4\,H_0$ (40\% of the scheduled headway), the bound Rodriguez et al.\ used to prevent excessive intervention. This cap matches the maximum value in the MARL action set $\Omega = \{0, 0.1, 0.2, 0.3, 0.4\}$, so EH and MARL are constrained to the same maximum hold and the comparison is fair. EH represents globally-aware reactive control: it uses information about both neighbors but follows a fixed rule rather than a learned policy. Because Eq.~\eqref{eq:eh} is a fixed linear rule computed from the currently observed gaps, EH is expected to handle breakdowns better than FH, since the enlarged backward gap $\hat{h}^+$ enters the rule directly, but it has no mechanism to anticipate or discount the heavier-tailed delays that the weather generator introduces, since the rule was not designed with a heavy-tailed disturbance in mind.
\clearpage
\subsubsection{Baseline Comparisons}

The three controllers span the comparison space. NC tests whether intervention helps at all. FH tests whether MARL outperforms classical local reactive control. EH tests whether MARL outperforms classical global reactive control and EH is the stronger of the two classical baselines, consistently reported as the best non-RL controller in the literature \cite{Wangsun,Rodriguez2023Cooperative}. The acceptance criterion for the MARL policy in Stage A (Section~\ref{subsec:evaluation}) is therefore stated relative to EH: matching or beating EH is the meaningful test.

This three-way comparison also allows the study to attribute MARL's advantage. If MARL beats NC but not FH or EH, the benefit is from intervening, not from learning. If MARL beats FH but not EH, the benefit is from global awareness, not from learning. Only if MARL beats EH does the benefit come from the learned policy itself. The same logic applies in Stage B: the disturbance-handling claim is meaningful only if MARL's advantage over EH grows with $\eta$, since both controllers see the same disturbance realizations.

State-of-the-art MARL methods such as CAAC and IQNC-M \cite{Wangsun} were considered as additional baselines but excluded. Including them would shift the research question from ``does MARL help on EDSA?'' to ``which MARL variant is best?'', a different question requiring different experimental data. 

\subsection{Data Analysis Methods}

The three response variables logged per run are defined as follows. \textbf{Mean passenger waiting time} ($\bar{W}$) is the average time elapsed from a passenger's arrival at a stop to their successful boarding, averaged across all passengers served and all stops over one simulated operating day:

\begin{equation}
\bar{W} = \frac{1}{P} \sum_{p=1}^{P} \left(t_p^{\text{board}} - t_p^{\text{arrive}}\right)
\label{eq:waiting_time}
\end{equation}

where $P$ is the total number of passengers served in the run and $t_p^{\text{board}}$, $t_p^{\text{arrive}}$ are the boarding and arrival times of passenger $p$. \textbf{Mean total travel time} ($\bar{T}$) is the average elapsed time from a bus's departure from the origin terminal to its arrival at the final stop of the sub-corridor, averaged across all bus trips completed during the simulated day. \textbf{Headway coefficient of variation} ($CV_h$) measures headway regularity:

\begin{equation}
CV_h = \frac{\sigma_h}{\mu_h}
\label{eq:headway_cv}
\end{equation}

where $\sigma_h$ and $\mu_h$ are the standard deviation and mean of observed inter-bus headways recorded at all stops over one simulated operating day. $CV_h = 0$ denotes perfectly regular headways; larger values indicate increasing bunching severity. This construction mirrors the baseline coefficient of variation $CV_0$ already defined for travel time (Table~\ref{tab:notation}), applied here to the headway distribution instead.

For each (control strategy, disturbance level) cell, $N \geq 30$ independent Monte Carlo runs are executed using matched random seeds across strategies. Three response variables are logged per run: mean passenger waiting time, mean total travel time, and headway coefficient of variation. The choice $N \geq 30$ ensures that the Central Limit Theorem applies to the run-mean of each response variable \cite{Wasserman2004}.

Statistical significance of performance differences between controllers will be assessed at a significance level of $\alpha = 0.05$. Because the same random seeds are reused across controllers within a disturbance level, comparisons between two controllers at the same disturbance level are naturally paired: the only difference between two runs is the controller, not the disturbance realization. Paired-sample tests will therefore be used for within-disturbance-level comparisons; unpaired tests will be used where seed matching does not apply (for example, comparing the same controller across different disturbance levels). Multiple-comparison correction will be applied across the four controllers to control the family-wise error rate.

For non-ideal disturbance levels, performance degradation is computed as the percentage change in each response variable relative to the same controller's $\eta = 0$ ideal-condition result. Confidence intervals on this degradation ratio will be computed by bootstrap resampling, since ratios of sample means are not guaranteed to be normally distributed and may be skewed under high-disturbance conditions \cite{EfronTibshirani1993}. The specific hypothesis tests, normality checks, multiple-comparison procedure, and bootstrap parameters will be selected during the implementation phase in consultation with the research supervisor, following recent guidance on reinforcement-learning evaluation~\cite{Henderson2018Matters,Agarwal2021Precipice}. Analysis is performed in Python using NumPy, SciPy, and pandas; raw run logs and analysis notebooks are version-controlled for reproducibility.

\subsection{Evaluation Methods}
\label{subsec:evaluation}



\textbf{Stage A: Ideal-condition evaluation.} The weather and breakdown generators are disabled; demand and traffic-speed scaling are fixed at baseline means ($\eta = 0$). The trained MARL policy and the three baselines are each run for $N \geq 30$ Monte Carlo iterations with matched seeds, reported in the format shown in Figure~\ref{fig:stage-a-illustrative}. Under ideal conditions EH is already near-optimal, so parity with EH is an acceptable Stage A outcome. Outperforming EH is the target but not required to pass Stage A, since the main contribution of this study is the characterization of performance under disturbance in Stage B.

\paragraph{Stage A acceptance criterion}
\begin{itemize}
\item[(i)] Mean passenger waiting time no worse than EH (no statistically significant degradation at $p < 0.05$ with multiple-comparison correction), and ideally a statistically significant improvement.
\item[(ii)] A statistically significant reduction in headway coefficient of variation relative to NC.
\end{itemize}

\begin{figure}[htbp]
    \centering
    \includegraphics[width=0.8\textwidth]{Figures/meth_fig_eo3_1_ideal_results.pdf}
    \caption[Illustrative Stage A output format]{Illustrative Output of the Stage A
    ideal-condition evaluation output, comparing the four controllers (NC, FH, EH,
    MARL) across the three response variables: (a) mean passenger waiting time,
    (b) mean total travel time, and (c) headway coefficient of variation. Error
    bars represent 95\% confidence intervals over $N \geq 30$ Monte Carlo runs
    with matched seeds. Numerical values shown are illustrative and do not
    reflect measured results. Authors' illustration.}
    \label{fig:stage-a-illustrative}
\end{figure}

\textbf{Stage B: Non-ideal-condition evaluation.} The trained MARL policy and the three baselines are evaluated across the swept disturbance range $\eta \in \{0.3, 0.6, 1.0, 1.3\}$, with the breakdown generator active at each level, using the Monte Carlo evaluation procedure illustrated in Figure~\ref{fig:aec-eval}. For each $(\eta, \text{controller})$ cell, $N \geq 30$ Monte Carlo iterations are run with matched seeds. Reported outputs are: (a) mean waiting time, travel time, and headway CV as functions of $\eta$ for each controller; (b) percentage degradation at each $\eta$ relative to the controller's own ideal-condition result; (c) the disturbance level at which MARL's advantage over the best-performing baseline becomes statistically significant; and (d) the degradation slope, defined as the rate of waiting-time increase per unit increase in $\eta$, with smaller slopes indicating greater resilience. An ablation is also run with each disturbance enabled alone (demand-only, weather-only, breakdown-only) at $\eta = 1.0$ to attribute degradation to disturbance type.

\paragraph{Stage B acceptance criterion} The MARL policy is reported as performing well under disturbance if its mean waiting time stays below that of the best-performing baseline across the full sweep, with statistical significance under the same correction as Stage A.

\begin{figure}[htbp]
    \centering
    \includegraphics[width=1\textwidth]{Figures/fig_3_2b_aec_evaluation.pdf}
    \caption[Monte Carlo evaluation layer]{The Monte Carlo evaluation layer. The frozen policy $\theta$ selects the highest-valued action at each control stop without exploration or weight updates. Matched seeds ensure that every controller faces the same disturbance realization, so performance differences reflect the control strategy alone. Three response variables are recorded per run: mean passenger waiting time, mean total travel time, and headway coefficient of variation. Authors' illustration.}
    \label{fig:aec-eval}
\end{figure}


\section{Methodological Challenges}
\label{sec:challenges}

Three challenges shape the experimental design. The design is scoped to remain feasible on accessible hardware (a single cloud-notebook GPU runtime), and a tiered scope plan (stated at the end of this section) guarantees a defensible minimum deliverable even if the full sweep cannot be completed.

\textbf{Q-value overestimation under heavy-tailed returns.} Value-based methods such as DQN overestimate Q-values \cite{vanHasselt2016DDQN}, and this bias grows when returns are heavy-tailed, as they are under high-intensity disturbance \cite{Henderson2018Matters,Agarwal2021Precipice}. The DDQN double estimator addresses this by decoupling action selection from evaluation \cite{vanHasselt2016DDQN}. Because disturbances are active during training, reward clipping and target-network soft updates are used as additional stabilizers throughout, not only at the highest intensities.

\textbf{Computational cost.} The evaluation budget follows directly from the experimental factors. The full $\eta$ grid spans five values $\{0, 0.3, 0.6, 1.0, 1.3\}$, with $\eta = 0$ serving as the Stage A ideal condition and the remaining four levels constituting the Stage B non-ideal sweep. Stage A runs four controllers at $N \geq 30$ each, giving $4 \times 30 = 120$ runs. Stage B runs four controllers across four non-ideal levels at $N \geq 30$ per cell, giving $4 \times 4 \times 30 = 480$ runs. The single-disturbance ablation (Section~3.2.10) adds three conditions (demand-only, weather-only, breakdown-only) across the same four controllers at $N \geq 30$, giving $3 \times 4 \times 30 = 360$ runs, for a total of approximately 960 evaluation runs. Two decisions keep this tractable. First, SUMO is restricted to one-time calibration; all Monte Carlo runs execute in the lightweight Python environment. Second, evaluation runs are CPU-bound and parallelizable across random seeds, so training is the only phase requiring sustained GPU time.

\textbf{Timeline feasibility.} The full evaluation suite of approximately $960$ runs is expected to require on the order of single-digit hours of single-core CPU time, given the event-driven simulation rates reported in comparable work \cite{Wangsun,Rodriguez2023Cooperative}; parallelizing across seeds on a multi-core runtime fits the full sweep within a single working session. MARL training is sized to the episode budget reported by Rodriguez et al.~\cite{Rodriguez2023Cooperative}, who train for $500 + 300 = 800$ episodes in two phases. The total training budget, including hyperparameter tuning and reward-weight iteration (EO 2.1), is projected at approximately 30 hours of T4-class GPU time, which fits within standard cloud-notebook allocations without institutional HPC resources. These time figures are planning estimates to be confirmed empirically during implementation rather than measured values. A tiered scope plan is maintained: the minimum deliverable is Stage A ($\eta = 0$) plus a two-level non-ideal sweep at $\eta \in \{0.6, 1.0\}$; the full four-level non-ideal sweep and ablation series are added if calibration and training proceed on schedule.


%\chapter{Results}
\hspace{\parindent} In this chapter, the author presents the results and detailed discussion of the performance evaluation for individual modules of the proposed \textit{NeuroSEE} pipeline. A more summarised discussion of these results can be found in the succeeding chapter.

\section{Motion Artefact Correction}
\hspace{\parindent} The main purpose of correcting movement artefacts is to aid the detection and segmentation of cells, which, in turn, can facilitate the accurate extraction of neural activity information for downstream analysis. Motion artefacts are usually induced when the mouse is turning or switching directions. Without prior correction, the succeeding stages in the pipeline may not work as accurately as expected. This, however, must be made as fast and computationally efficient as possible, especially when large high-dimensional datasets are being processed and analysed. 

For the purpose of comparison, the author attempted to implement a Hidden Markov Model-based line-by-line correction using the method presented in \cite{dombeck2007imaging}. However, this method was shown to be inapplicable to large high-resolution datasets, as it is very computationally expensive and takes a large amount of time to finish correcting tens of thousands of image frames. Aside from this issue, it also poses a problem caused inherently by raster scanning. After correction, it was found that some frames contain "blank" lines that may represent the lines that were not sampled on that particular frame during the instance of animal movement, i.e. the animal moved in such a way that the laser scan hopped over that region. Since the pixels are acquired serially during raster scanning, there may be no actual data recorded from those lines. Hence, there may be no other way to reconstruct them, thereby producing such blank lines. This, however, may be addressed by performing translational registration using frame-wise phase correlation prior to performing the line-by-line correction. However, it will make the process much more complex and will require more computational power and processing time, making the method inefficient for large scale datasets. 

Thus, the author has successfully shown that using a subpixel image registration, while taking advantage of both the structural marking (red) and functional imaging (green) channels, can lead to a fast and accurate motion artefact correction, making it suitable for larger datasets. 

\subsection{\textit{NeuroSEE} can accurately correct movement artefacts in \textit{in vivo} imaging datasets from awake, head-fixed moving animals}
\hspace{\parindent} To evaluate whether the proposed \textit{NeuroSEE} pipeline can correct movement artefacts, the author tested it on \textit{in vivo} two-photon calcium imaging data which were recorded while imaging an awake, head-fixed moving mice. 

The proposed method operates by splitting each frame into overlapping patches. The author initially sets the patch size to $128 \times 128$ pixel with an additional 32 pixels of overlap in each direction. Each of the patches were registered against a template. The initial template was computed by taking the median of $n=200$ frames in the structural marking (red) channel. After registering the patches, each path is further split into smaller overlapping patches. The patch sizes were modified experimentally to identify the size that gives the best performance. It has been shown that larger initial patches (either $128 \times 128$ or $96 \times 96$) are a good choice, especially when the SNR of the dataset is not high enough. 

The performance of this method can be shown in Figure ~\ref{fig:Res1}. Panels A and B show the mean images of the red channel before and after correction, respectively. Uncorrected data gives a blurred mean image, while a properly motion-corrected data results to a crisp, noiseless mean image that indicates stability. To evaluate the improvement in the stability of the imaging data after correction, the author computed the correlation of each of the image frames with respect to the mean image. It can be shown in Figure ~\ref{fig:Res1}C that there is a significant increase in (\ac{cwm}) between the motion-corrected and raw data. Intuitively, this increase in \ac{cwm} coefficient indicates an improvement in stability. Figure ~\ref{fig:Res1}D and ~\ref{fig:Res1}E provide a more detailed frame-by-frame correlation with the mean. As shown in ~\ref{fig:Res1}D, all of the frames in the motion-corrected data had an increase in \ac{cwm}. Each sample point in this plot shows the correlation between each frame and their corresponding mean image. Points lying above the red line indicate an increase in \ac{cwm}.

\begin{figure*}[!t]
	\centering
	\includegraphics[width=5.7in]{figures/figs/Res1.png}
	\caption{Performance of \textit{NeuroSEE} in correcting motion artefacts}
	\label{fig:Res1}
	\captionsetup{singlelinecheck=off,font=small}
	% EDIT THIS:
	\caption*{\textit{NeuroSEE} stabilises image sequences that were corrupted by motion artefacts: (A) Mean image of the structural marking (red) channel of the raw data; (B) A sharper and stabilised mean image of the structural marking (red) channel of the motion-corrected data; (C) Correlation coefficients of the raw and motion-corrected frames with respect to their corresponding means; (D-E) Frame-by-frame correlation-with-mean of the raw and motion-corrected data}
\end{figure*}

\subsection{\textit{NeuroSEE}'s movement correction leads to better source extraction}
\hspace{\parindent} Most of the existing methods for source extraction assume stable image sequences, i.e. each cell retains its position throughout all the frames in the video. To examine the effect of motion artefact correction to cell detection and segmentation, the author applied the proposed hybrid detection algorithm to raw and motion-corrected data. It has been shown that raw datasets can lead to cells not being detected and segmented because the displacement of the position of cells across frames may lead to a blurred mean and correlation images, which can then result to improper evolution of level set functions. The quality of the extracted calcium traces significantly increase after motion artefact correction. 

Figure ~\ref{fig:Res3}A shows the mean image of the red channel with overlaid contours nine representative cells. Their respective extracted ratiometric calcium traces are shown in Figure ~\ref{fig:Res3}B. It can be shown that \textit{NeuroSEE} can accurately detect and segment different types of cells and extract their corresponding calcium transients with sufficient signal-to-noise ratio. 

\section{Cell Detection and Segmentation}
\hspace{\parindent} The goal of proposing the \textit{NeuroSEE} pipeline is to standardise the analysis to a wide variety of cell models and imaging data types. To address this, the pipeline must not heavily rely on cellular morphology nor in stereotypical temporal activity. Hence, the author integrated morphology- and activity-based approaches to create a more flexible and versatile analysis pipeline.

The performance of the proposed algorithm was optimised by experimentally tuning its parameters. Alpha, $\alpha$, parameter is sensitive to the degree of correlation of the detected cells. The higher the value of $\alpha$, i.e. approaching 1, the more conservative (or more selective) the system is in selecting a candidate ROI. This means that only those with high correlation across time will be nominated as candidate ROIs. On the other hand, if $\alpha$ is set to a lower value, i.e. approaching 0, more cells will be detected from the correlation image. Based on the experiments, it was found that a value of $\alpha$ in the range of 0.3 - 0.6 gives better initialisation results, as lower $\alpha$ values overestimate the number of ROIs, which can then be automatically pruned during the evolution of the level set functions. 

Lambda, $\lambda$, parameter, on the other hand, highly depends on the cell’s temporal activity. The higher the value of $\lambda$ the more weight the data-based term has in the cost function, resulting to a more defined active contour. Based on the experiments, it was found that a value of $\lambda$ in the range of 10 - 60 gives better final cell boundary estimates.

During testing, \textit{NeuroSEE} detected a total of 321 ROIs from the manually-labelled imaging dataset of the Neurofinder Challenge (Figure ~\ref{fig:Res2}C). Automatic initialisation using the morphology-based approach produced 120 ROIs (Figure ~\ref{fig:Res2}A), while activity-based approach produced produced 248 ROIs (Figure ~\ref{fig:Res2}B), summing to 368 initial ROI estimates with 47 automatically pruned or merged during the level set segmentation. Figure ~\ref{fig:Res2}D shows the manually labelled cells. 

\subsection{\textit{NeuroSEE} outperforms existing approaches}
\hspace{\parindent} The performance of \textit{NeuroSEE}'s ROI segmentation module was compared with other existing approaches, such as ABLE \cite{reynolds2017able}, Suite2p \cite{pachitariu2016suite2p}, and CNMF \cite{pnevmatikakis2016simultaneous}. As shown in Figure ~\ref{fig:Res2}E and Table ~\ref{table:Table1}, \textit{NeuroSEE} achieved the highest success rate, precision, and recall values among all these methods. These were computed by comparing their results to the manual labels provided in the Neurofinder Challenge dataset. It has been shown clearly that \textit{NeuroSEE} outperforms the other approaches, as it can detect even the non-circular/non-elliptical cells, as well as the temporarily inactive cells. Hence, \textit{NeuroSEE} is robust to cell shape and to temporal cellular inactivity. 


\subsection{\textit{NeuroSEE} eliminates neuropil contamination}
\hspace{\parindent} \textit{NeuroSEE} is capable of eliminating neuropil contamination prior to the extraction of the calcium traces. Figure ~\ref{fig:Res3}C shows a representative cell (cell 1 in Figure ~\ref{fig:Res3}A and B). The neuropil region, i.e. the region outside the active contour bounded by the yellow dashed line, is divided into four subregions. The signals coming from these subregions are weighted. This is considered as the neuropil signal (right-middle; red). This signal is subtracted from the measured signal of the ROI (right-top; magenta). The resulting signal (right-bottom; blue) is the decontaminated signal. It can be observed that the decontamination of the extracted calcium traces gives a higher signal-to-noise ratio. 

\begin{sidewaysfigure*}
	\centering
	\includegraphics[width=7in]{figures/figs/Res2.png}
	\caption{Performance of \textit{NeuroSEE} in detecting and segmenting cells}
	\label{fig:Res2}
	\captionsetup{singlelinecheck=off,font=small}
	% EDIT THIS:
	\caption*{\textit{NeuroSEE} gives the best performance among the other three existing cell detection and segmentation algorithms: (A) Detected cells during morphology-based initialisation; (B) Detected cells during activity-based initialisation; (C) Detected cells after performing the proposed hybrid cell detection and segmentation method; (D) Manually-labelled cells from the Neurofinder Challenge dataset; (E) Comparison of the performance metrics for different algorithms}  
\end{sidewaysfigure*}

\begin{figure*}[!t]
	\centering
	\includegraphics[width=6in]{figures/figs/Res3.png}
	\caption{Performance of \textit{NeuroSEE} in extracting neural activity information}
	\label{fig:Res3}
	\captionsetup{singlelinecheck=off,font=small}
	% EDIT THIS:
	\caption*{\textit{NeuroSEE} accurately extracts calcium transients: (A) Summary image of a two-photon video showing representative cells that are detected by the proposed algorithm; (B) Ratiometric calcium traces extracted from the nine cells indicated in the left image; (C) A $100 \times 100$ pixel patch showing one cell and its neuropil subregions. The right panel shows the composite signal prior to neuropil decontamination (top; magenta), the neuropil signal (middle; red), and the decontaminated signal (bottom; blue).}
\end{figure*}

\begin{table}[t]
\centering
\begin{tabular}{ |p{2.5cm}||p{3.5cm}|p{3cm}|p{3cm}| }
 	\hline
	 \multicolumn{4}{|c|}{Performance Evaluation} \\
 	\hline
 	Method &Success Rate (\%) &Precision (\%)&Recall(\%)\\
	 \hline
 	NeuroSEE   & 75.4 & 72.6 & 78.4\\
 	ABLE   & 67.5 & 67.5 & 67.5\\
 	CNMF   & 63.4 & 60.7 & 66.5\\
 	Suite2p   & 63.7 & 56.5 & 73.1\\
	\hline
\end{tabular}
\caption{Algorithm success rate on manually labelled dataset}
\label{table:Table1}
\end{table}




\section{Place Field Analysis}
\hspace{\parindent} To further evaluate whether the proposed \textit{NeuroSEE} pipeline is capable of extracting accurate neural activity information from the detected cells, the author performed 2D histogram-based place field analysis. This type of histogram binning-based technique is highly dependent on the amount of firings during the recording. Hence, it is necessary that the animal under test is active during the imaging sessions, and that it traverses through almost each of the regions in the environment so as to get a good map-like representation of space.  

In order to perform place field analysis, it is required to simultaneously track the path of the mouse as it explores the circular chamber during the two-photon imaging sessions. This is done using a magnetic tracker whose sampling rate is 100 Hz. Since the frame scanning rate and the sampling rate of the tracker are not the same, the first thing to do before performing place field mapping is to make sure that their time stamps are synchronised. This is done by the \textit{Processor} module at the \textit{Behavioural Tracking} stage of the \textit{\ac{neurosee}} pipeline. When the \ac{2p} spike onset time series and the tracking time series are in sync, the occupancy, rate and place maps can now be constructed. The estimation of place fields can be used to quantify information encoding and transmission. 

The sizes of the occupancy, rate and place maps were experimentally modified during the analysis. It was found that a $100 \times 100$ bin map size gives sufficient results. Each bin in the occupancy and rate maps gives the frequency of space occupancy and cell firing, respectively. Dividing the rate map by the occupancy map produces the place map, which suggests the selectivity of firing of individual place cells with respect to a particular region in space. 

\subsection{\textit{NeuroSEE} shows promising performance in monitoring the dynamics of hippocampal place cells}
\hspace{\parindent} Based on the place maps obtained after performing place field analysis, it has been shown that \textit{NeuroSEE} can be used to monitor the changes in the spiking activity of hippocampal place cells across several imaging sessions. The place maps show promising results in demonstrating the spatial selectivity of a population of cells in the hippocampus. The estimation of place fields can be used in quantifying information encoding and transmission. Hence, this tool can be used in investigating how the spatiotemporal dynamics of neural cortical circuits involve in memory encoding and recall are affected by Alzheimer's disease as it progresses. Place field analysis can be performed across several imaging days to quantify and analyse the possible changes in dynamics of a population of place cells. 

To further improve the results, it is suggested to record more imaging data for each imaging session. The longer the time of recording, i.e. the longer the mouse explores the environment, the more pronounced the place maps would be. It is also suggested to use a more accurate and precise estimation of the spatial information content by using information theoretic-based approaches, such as Gaussian Process-based method. This may provide a more data-efficient estimation than 2D spike binning-based approach. 

\begin{figure*}[!t]
	\centering
	\includegraphics[width=5.7in]{figures/figs/Res4.png}
	\caption{Performance of \textit{NeuroSEE} in place field analysis}
	\label{fig:Res4}
	\captionsetup{singlelinecheck=off,font=small}
	% EDIT THIS:
	\caption*{(A) Behavioural track: Gray lines represent the trajectories of the mouse in the circular chamber and the different coloured dots represent the points in the chamber where different respective cells fired action potentials; (B) Occupancy map: Indicates the frequency of the mouse in occupying a certain region in space. The space is divided into spatial bins.; (C) Spiking activities of representative cells: The signal on top is the extracted ratiometric calcium trace and the signal below is the non-negative derivative of the calcium trace which represents the spike onsets.; (D-E) Original and smoothened place maps}
\end{figure*}

%\chapter{Discussion and Conclusions}

\hspace{\parindent} In this research project, the author presents a standardised computational pipeline that is capable of accurately extracting activity information from automatically segmented cells in performing a comprehensive analysis of two-photon calcium imaging data of any cell model. The proposed pipeline addresses the issues of movement artefacts during imaging of awake mobile mice, the addition of neuropil contamination due to out-of-focus fluorescence from the surrounding neuropil and neighbouring cells, and the heterogeneity in the cellular properties, morphology and temporal activity of the regions of interest. All of these issues impede the accurate detection and segmentation of cells, leading to inaccurate extraction of their neural activity information. 

Among these issues, the problem of movement artefacts poses the greatest challenge in developing an accurate automated analysis pipeline, as the accuracy of the succeeding stages highly depend on the stability of the image sequence, i.e. the positions of the regions of interest do not change over time. If this issue is not properly addressed, the cells will not be properly detected either using morphology-based or activity-based segmentation approaches, or even both. Detecting the cell boundaries using either of these approaches rely on the mean and correlation images of image sequences. Unstable image sequences, i.e. uncorrected data corrupted with either low or high amount of movement artefacts, result to a blurred and noisy mean and correlation images. The evolution of level set functions during the update phase of the level set method is also affected by the instability of the positions of the ROIs of uncorrected data. 

In this study, the author addressed the movement artefacts by applying a subpixel image registration technique that is capable of correcting not only rigid artefacts caused by the natural uniform motion of the brain, but more importantly, it is capable of correcting non-rigid artefacts that arise from the finite speed of raster scanning and the non-rigid deformations due to mechanical strain and head displacement during \textit{in vivo} imaging of head-fixed moving mice. Unlike the HMM-based line-by-line correction approach, \textit{NeuroSEE}'s motion correction algorithm is faster and computationally efficient - making it a suitable method for large high-resolution, high-dimensional imaging datasets. This has been successfully demonstrated when this method was applied to a motion artefact-induced \textit{in vivo} imaging dataset from an awake head-fixed moving mice. The motion-corrected data produced a crisp, more focused and noiseless mean image, indicating stability across the frames in the videos. A significant increase in the correlation-with-mean metric of the motion-corrected data was also shown when compared to that of the raw data. This significant increase validates the significant improvement in the stability of the resulting corrected data. 

In order to validate this further, the proposed hybrid cell detection method in the \textit{NeuroSEE} pipeline was implemented to both raw and motion-corrected data and compared their results. It was shown that the performance of the ROI segmentation module in the motion-corrected data outperforms that of the raw data - proving how the motion correction module improves the detection of cells, as well as the extraction of their neural activity information.

To generalise and make the cell detection more flexible and versatile to a wide variety of imaging datasets, the author integrated both morphology- and activity-based approaches. This is based on the intuition that combining these approaches can improve the general appliabillity of ROI segmentation tools by taking advantage of their individual capabilities, i.e. by detecting temporarily inactive cells using the morphology-based approach and detecting non-circular or non-elliptical-looking cells using the activity-based approach. The integration of these approaches requires the use of both the functional imaging (green) and structural marking (red) channels, a property of the proposed pipeline that is not commonly seen in existing tools. 

Accurate cell detection requires proper tuning of parameters. In this study, the $\alpha$ parameter was tuned optimally to define the appropriate selectivity of initialisation. This parameter is tuned in such a way that it is not too high, i.e. too selective and conservative, resulting to some cells not being detected; but at the same time, it must not be set too low such that it causes overestimation of ROIs, which can make the evolution of the active contours much more complicated and time-consuming. 

After detecting the cells using the hybrid approach, the initial cell boundary estimates, which are represented by active contours, are updated by evolving the level set functions. The criteria in evolving these level set functions are based both on the morphology- and activity-based properties of the two imaging channels. The $\lambda$ parameters was also tuned in order to give a better, smoother final cell boundary estimates. 

In order to evaluate the performance of the proposed \textit{NeuroSEE}'s ROI segmentation module, it was tested on a manually labelled dataset containing manual labels of the expected ROIs which can serve as the ground truth. The predicted cells are matched with the ground truth cells to compute its performance metrics. It was then compared to existing activity-based and morphology-based approaches. It has been shown that \textit{NeuroSEE} has clearly outperformed the other methods in detecting and segmenting cells. It obtained a success rate of 75.4\%, a precision of 72.6\%, and a recall value of 78.4\%, as compared to the other methods whose success rates are all less than 70\%. 

\textit{NeuroSEE} has also successfully decontaminated the extracted calcium traces from neuropil artefacts. By subtracting the weighted subregion neuropil signals from the measured calcium signals, By doing so, the signal-to-noise ratio of the decontaminated signal has improved, leading to a more accurate downstream analysis, such as place field mapping.

Finally, in order to validate whether the extracted calcium signals can be used to analyse the dynamics of the place cells in the hippocampus during mouse exploration and to demonstrate whether these cells in the CA1 region can provide a map-like representation of space, a 2D histogram-based place field mapping was performed. The non-negative derivative of the extracted calcium traces were taken to estimate the onset of spiking of action potentials. These spike onsets were used to construct the rate map of individual cells, which when divided by the occupancy map produces their respective place field maps. With the place maps obtained from the collected neural activity information, it can be shown that \textit{NeuroSEE} is capable of providing an accurate analysis of the dynamics of a population of cells that were detected and segmented automatically. Place field estimation can be used to further quantify information encoding and transmission. 

Thus, the proposed pipeline has shown promising results and has successfully demonstrated that efficient artefact correction and the use of a hybrid approach in detecting and segmenting cells lead to a more accurate and generally applicable extraction of neural activity information. 
%\chapter{Future Work}

\hspace{\parindent} In the future, the author aims to integrate within the current \textit{NeuroSEE} pipeline a novel information theoretic and machine learning-based algorithm for reducing the dimensionality of the large-scale multivariate signals extracted, as well as an algorithm that is capable of estimating information flow graphs that reflect disease progression. 

For the place field analysis, the author aims to further improve the algorithm by using information theoretic approaches in accurately estimating the spatial information content, e.g. using Gaussian Process-based method. It would also be helpful if other information theoretic analysis are integrated within the pipeline, e.g. spike inference and spike entropy calculation, in order to quantify information encoding and transmission.

On a more general note, the author aims to use the proposed \textit{NeuroSEE} pipeline in investigating how the spatiotemporal dynamics of neural cortical circuits involved in memory encoding and recall are affected by Alzheimer's disease as it progresses. More specifically, it would be important to characterise the differences between different mouse models of AD in terms of their effects on information flow in hippocampal-cortical circuits, so as to generalise our understanding on the effect of the different manifestations of AD on memory circuit dynamics. 

In mouse models of AD, amyloid or tau protein-related pathologies result in changes in cortical circuitry, which are particularly evident in the hippocampus. What is not known is how these changes relate to the function of this circuitry in memory encoding and recall processes, or how therapies which have been found to restore behavioural performance in mouse models affect this circuitry. Hence, it is good to examine the effects of therapeutic interventions, including pharmacological agents and neuromodulatory therapies, on the information processing properties in different mouse models. 

All of the above-mentioned works will be performed as the author progresses towards the PhD phase of the Centre for Doctoral Training in Neurotechnology for Life and Health programme. 
%\appendix
\chapter{Other Supplementary Figures}

\chapter{Source Codes}

%\clearpage
\printglossary[type=\acronymtype]


\bibliography{thesis_refs}
\bibliographystyle{ieeetr}


\end{document}